% Options for packages loaded elsewhere
\PassOptionsToPackage{unicode}{hyperref}
\PassOptionsToPackage{hyphens}{url}
\PassOptionsToPackage{dvipsnames,svgnames,x11names}{xcolor}
%
\documentclass[
  11pt,
]{article}
\usepackage{amsmath,amssymb}
\usepackage{iftex}
\ifPDFTeX
  \usepackage[T1]{fontenc}
  \usepackage[utf8]{inputenc}
  \usepackage{textcomp} % provide euro and other symbols
\else % if luatex or xetex
  \usepackage{unicode-math} % this also loads fontspec
  \defaultfontfeatures{Scale=MatchLowercase}
  \defaultfontfeatures[\rmfamily]{Ligatures=TeX,Scale=1}
\fi
\usepackage{lmodern}
\ifPDFTeX\else
  % xetex/luatex font selection
\fi
% Use upquote if available, for straight quotes in verbatim environments
\IfFileExists{upquote.sty}{\usepackage{upquote}}{}
\IfFileExists{microtype.sty}{% use microtype if available
  \usepackage[]{microtype}
  \UseMicrotypeSet[protrusion]{basicmath} % disable protrusion for tt fonts
}{}
\makeatletter
\@ifundefined{KOMAClassName}{% if non-KOMA class
  \IfFileExists{parskip.sty}{%
    \usepackage{parskip}
  }{% else
    \setlength{\parindent}{0pt}
    \setlength{\parskip}{6pt plus 2pt minus 1pt}}
}{% if KOMA class
  \KOMAoptions{parskip=half}}
\makeatother
\usepackage{xcolor}
\usepackage[margin=24mm]{geometry}
\setlength{\emergencystretch}{3em} % prevent overfull lines
\providecommand{\tightlist}{%
  \setlength{\itemsep}{0pt}\setlength{\parskip}{0pt}}
\setcounter{secnumdepth}{-\maxdimen} % remove section numbering
\ifLuaTeX
\usepackage[bidi=basic]{babel}
\else
\usepackage[bidi=default]{babel}
\fi
\babelprovide[main,import]{english}
% get rid of language-specific shorthands (see #6817):
\let\LanguageShortHands\languageshorthands
\def\languageshorthands#1{}
\usepackage{amsmath,amssymb,amsthm,mathtools,mathrsfs}
\usepackage{lmodern,microtype,xurl,fancyhdr,needspace}
\setlength{\emergencystretch}{3em}
\setlength{\parindent}{0pt}
\setlength{\parskip}{0.4em}
\setlength{\headheight}{14pt}
\allowdisplaybreaks[1]
\pagestyle{fancy}
\fancyhf{}
\fancyhead[L]{\small Stronger symmetric projection stability}
\fancyhead[R]{\small mxym}
\fancyfoot[C]{\thepage}
\renewcommand{\headrulewidth}{0.3pt}
\ifdefined\pdfinfoomitdate\pdfinfoomitdate=1\fi
\ifdefined\pdfsuppressptexinfo\pdfsuppressptexinfo=15\fi
\ifdefined\pdftrailerid\pdftrailerid{}\fi
\ifLuaTeX
  \usepackage{selnolig}  % disable illegal ligatures
\fi
\usepackage{bookmark}
\IfFileExists{xurl.sty}{\usepackage{xurl}}{} % add URL line breaks if available
\urlstyle{same}
\hypersetup{
  pdftitle={Stronger symmetric projection stability},
  pdfauthor={mxym},
  pdflang={en},
  colorlinks=true,
  linkcolor={Maroon},
  filecolor={Maroon},
  citecolor={Blue},
  urlcolor={Blue},
  pdfcreator={LaTeX via pandoc}}

\title{Stronger symmetric projection stability}
\author{mxym}
\date{7 October 2026}

\begin{document}
\maketitle

This complete proof collection establishes the two globally valid bounds
\(D(K,\mathcal E_d)-1\le\min\{d^{12}\delta^{1/(3d)},d^{19}\delta^{1/[3(d-1)]}\}\)
for origin-symmetric full-dimensional bodies in dimension \(d\ge3\). The
optimum remains unresolved. Section and equation labels below are local
to each named component. The first component states all assumptions and
constants; the second supplies the retained cone-law input and
whole-class obstruction. The historical exponent and gap statements in
Component II refer to the preserved baseline theorem; Component I gives
the stronger current modulus and narrower unresolved gap. No
quantitative inverse-Minkowski theorem is used.

\newpage

\section{I Main theorem and assembly}\label{i-main-theorem-and-assembly}

\subsection{A stronger upper-end exponent from weighted rank-two
matching}\label{a-stronger-upper-end-exponent-from-weighted-rank-two-matching}

7 October 2026. This note proves a strengthening of the completed
symmetric projection-cone stability theorem. The completed supplement is
unchanged. The new steps are a sharp square-root matching estimate and a
direct recovery theorem from off-block normals. No heuristic tuning of
thresholds is used.

\subsubsection{1. Statement}\label{statement}

Let \(K\subset\mathbb R^d\) be a centrally symmetric full-dimensional
convex body, centered at zero, and let \(a(K)\) be the entry005 affine
projection-cone invariant. Set \(\delta=1/2-a(K)\). Let \(\mathcal E_d\)
be the entire class of linear images of Cartesian products of centrally
symmetric convex factors of dimensions one and two. Use \[
d_{\rm BM}(K,E)=\inf\{\lambda\ge1:TE\subseteq K\subseteq\lambda TE,
\ T\in\mathrm{GL}(d)\},\qquad
D(K,\mathcal E_d)=\inf_{E\in\mathcal E_d}d_{\rm BM}(K,E).
\]

\textbf{Theorem 1 (weighted matching).} For every \(d\ge3\) and every
such \(K\), \[
\boxed{D(K,\mathcal E_d)-1
\le d^{12}\bigl(\tfrac12-a(K)\bigr)^{1/(3d)}.}                 \tag{1}
\] There is no smoothness, curvature, polytope, atom, or
absolute-continuity hypothesis. In dimensions one and two the distance
is one.

For the explicit unsimplified constants put \[
\begin{aligned}
N&=d+1,&\kappa_d&=\pi^{d/2}/\Gamma(1+d/2),&a_d&=d^{-2d},\\
Z_d&=2^{d+9}N\sqrt d\binom d3/\kappa_d,
&F_d&=13N+3Z_d,&B_d&=4NF_d/a_d,\\
\widehat L_d&=5d^2a_d^{-2}B_d^{1/2},
&\widehat H_d&=d(\sqrt d+1)\widehat L_d,
&\widehat C_d&=2d(2d\widehat H_d)^{1/d}.
\end{aligned}
\] The proof gives the stronger constant \(\widehat C_d\le d^{12}\), and
\[
D(K,\mathcal E_d)-1\le\widehat C_d\delta^{1/(3d)}.            \tag{2}
\] If \[
0<\delta\le\widehat\delta_d^*:=
 \left(\frac{2^{-d}}{2d\widehat H_d}\right)^3,
\] there is also the explicit containment estimate \[
D(K,\mathcal E_d)
\le\frac1{1-(2d\widehat H_d\delta^{1/3})^{1/d}}.              \tag{3}
\] The constants are conservative and depend on dimension. The earlier
all-class corner-truncation obstruction still excludes exponents above
\(1/d\). The next theorem further improves the geometric exponent.

\textbf{Theorem 2 (weighted matching and direct geometric recovery).}
Under exactly the same hypotheses, \[
\boxed{D(K,\mathcal E_d)-1
\le d^{19}\bigl(\tfrac12-a(K)\bigr)^{1/[3(d-1)]}.}             \tag{8}
\] Both estimates hold globally, so their minimum is also a valid
modulus. The remaining optimal-exponent interval is
\([1/[3(d-1)],1/d]\). The factor \(d^{19}\) is deliberately coarse.

\subsubsection{2. The inputs retained from the completed
proof}\label{the-inputs-retained-from-the-completed-proof}

The normalized cone law is \[
\nu_K=(u\mapsto u/h_K(u))_*
       \left[\frac{h_K(u)}{d|K|}\,dS_K(u)\right].
\] For independent samples \(X_i\) define \[
A=\mathbb E|\det(X_1,\ldots,X_d)|,\quad
B=\mathbb E\left|\det\begin{pmatrix}
X_1&\cdots&X_{d+1}\\1&\cdots&1
\end{pmatrix}\right|.
\] Entry005 v3 proves \(a(K)=B/((d+1)A)\) for general convex bodies. For
a \((d+1)\)-tuple, write \(\phi\) for its fourth-largest absolute
horizontal cofactor. The completed proof supplies the following three
established inputs.

\begin{enumerate}
\def\labelenumi{\arabic{enumi}.}
\tightlist
\item
  In symmetric John position \(B_2^d\subseteq K\subseteq\sqrt d B_2^d\),
  the law is supported in the unit ball and \(A\ge d^{-3d/2}\ge a_d\).
  This follows from the regular-contact identity
  \(\mathbb E[Y X^{\mathsf T}]=I/d\) and sequential wedge expectations.
\item
  The cone-specific contact-slab argument gives
  \(\mathbb E\phi\le F_d\delta^{2/3}\). Its conditional determinant
  denominator cancels before integration. The estimate includes singular
  tuples, nonatomic laws, flat-face atoms, and the zero-deficit case.
\item
  If \(Q\) is centered symmetric and full-dimensional,
  \(Q\subseteq RB_2^d\), and \(\mathcal W\) is the union of an
  orthogonal block decomposition, then its product of block projections
  \(P\) satisfies \[
  (1-(2dv)^{1/d})P\subseteq Q\subseteq P,
  \quad v=R(\sqrt m+1)\mathbb E_{\nu_Q}\operatorname{dist}(X,\mathcal W),
  \quad 2dv<1.                                             \tag{4}
  \] Here \(m\) is the number of blocks. The proof bounds
  \(V(Q[d-1],P)/|Q|-1\) by \(v\), uses Minkowski's first inequality, and
  inserts a homothetic image of a symmetric half of \(P\) into its
  missing cap. It assumes no inradius bound or inverse cone-volume
  uniqueness theorem.
\end{enumerate}

These inputs are Sections 2, 3 and 5 of the completed English proof. For
Theorem 1 the replaced argument is its Section 4 matching estimate;
Theorem 2 also replaces the final geometric recovery step. The full new
matching proof is in the accompanying \url{matching.md}, and its
essential quantitative mechanism is recorded next.

\subsubsection{3. Why weighted matching improves the
power}\label{why-weighted-matching-improves-the-power}

For a bounded Borel law with \(A\ge a>0\),
\(\mathbb E\phi\le\varepsilon\), put \(q=4(d+1)\varepsilon/a\). The
simultaneous Fubini-section selection in \url{matching.md} gives a
sampled basis \(U\), with \(D=|\det U|\ge a/2\), whose one-sample and
every omitted-basis two-sample section each have fourth-cofactor
expectation at most \(q\). In particular \[
\|U\|\le\sqrt d,\quad \|U^{-1}\|\le2\sqrt d/a,
\quad |\alpha_i|\le1/D,\quad
\mathbb Es_3(\alpha)\le q/D^2,\qquad\alpha=U^{-1}X,
\] where \(s_3\) is the third-largest coordinate magnitude.

Retain the two largest coordinates of each vector, with fixed-index tie
breaking. In an omitted-basis section the cofactors, up to signs, are \[
D\alpha_j,\quad D\beta_j,\quad
D(\alpha_i\beta_j-\alpha_j\beta_i)\quad(i\ne j).
\] Truncating both vectors changes each entry by at most
\(2s_3(\alpha)+2s_3(\beta)\). Absolute value and a fixed order statistic
are 1-Lipschitz in the maximum norm, so the truncated section
expectation is at most \(q+4q/D^2\).

Give a coordinate pair \(ij\) the weighted mass \[
m_{ij}=\mathbb E\left[
\min(|\alpha_i|,|\alpha_j|)
\mathbf1_{\{ij\text{ is the chosen top pair}\}}\right].
\] For adjacent pairs \(ij,jk\), their two truncated vectors make four
cofactors at least \(D^2r_\alpha r_\beta\), where \(r\) denotes the
smaller retained magnitude. Thus \[
m_{ij}m_{jk}\le\frac{q+4q/D^2}{D^2}\le\frac{5q}{D^4}.       \tag{5}
\] Pairs with mass exceeding \(\sqrt{q+4q/D^2}/D\) form a matching. Use
them as plane blocks and the remaining indices as line blocks. A
discarded pair contributes its smaller coordinate to distance from the
chosen block union, rather than its entire probability mass.
Consequently, for \(0\le q\le1\), \[
\begin{aligned}
\mathbb E\operatorname{dist}(\alpha,\mathcal W)
&\le\frac{\sqrt d\,q}{D^2}
 +\binom d2\frac{\sqrt{q+4q/D^2}}D\\
&\le\frac{\sqrt d+\sqrt5\binom d2}{D^2}\sqrt q
\le5d^2a^{-2}\sqrt q.                                     
\end{aligned}
\tag{6}
\] The last inequality follows from
\(4(\sqrt d+\sqrt5\binom d2)\le4d+5d(d-1)\le5d^2\). For \(q=0\), all
positive-weight edges form a matching and the third coordinate vanishes
almost surely. This proves the zero case without a limit. All selections
are finite Borel partitions, so no point needs positive mass.

\subsubsection{4. Complete assembly of the stronger
theorem}\label{complete-assembly-of-the-stronger-theorem}

Put \(K\) in symmetric John position, and apply (6) with \(a=a_d\),
\(\varepsilon=F_d\delta^{2/3}\). Then \[
q=B_d\delta^{2/3},\qquad
\mathbb E\operatorname{dist}(U^{-1}X,\mathcal W)
\le\widehat L_d\delta^{1/3}.                               \tag{7}
\] For \(\delta\le\widehat\delta_d^*\), the condition \(q\le1\) is
valid: \(\widehat H_d\ge B_d^{1/2}\), and therefore \[
q\le B_d\left(\frac{2^{-d}}{2d\widehat H_d}\right)^2\le1.
\]

Apply the primal map \(U^{\mathsf T}\). The body \(Q=U^{\mathsf T}K\) is
contained in \(dB_2^d\); its cone law is exactly \((U^{-1})_*\nu_K\).
Its coordinate line and plane blocks are orthogonal. Thus (4), with
\(R=d\), \(m\le d\), and (7), gives (3). Every projection factor is
centrally symmetric and full-dimensional in its block, so
\(P\in\mathcal E_d\). No decomposition has to be canonical or unique.

At the smallness threshold, \[
r:=(2d\widehat H_d\delta^{1/3})^{1/d}\le1/2.
\] The sandwich \((1-r)P\subseteq Q\subseteq P\) gives \[
D(K,\mathcal E_d)-1\le\frac r{1-r}
\le2(2d\widehat H_d)^{1/d}\delta^{1/(3d)}.
\] For every larger deficit use the John sandwich with a cube, which
gives \(D(K,\mathcal E_d)\le d\). The global coefficient in (2) covers
this range because \[
\widehat C_d(\widehat\delta_d^*)^{1/(3d)}=d.
\] At \(\delta=0\), (6) gives exact block support and (4) gives the
product itself. For \(d=1,2\), \(K\) is already an allowed factor. This
proves (2) and (3) in all cases.

To prove the stated numerical constant, the inscribed cube bound
\(\kappa_d\ge(2/\sqrt d)^d\) yields \[
Z_d\le2^{10}d^{d/2+9/2},\quad
F_d\le2^{12}d^{d/2+9/2},\quad
B_d\le2^{15}d^{(5/2)d+11/2}.
\] Consequently \[
\widehat L_d\le5\,2^{15/2}d^{(21/4)d+19/4},\qquad
\widehat H_d\le10\,2^{15/2}d^{(21/4)d+25/4},
\] and \[
\widehat C_d^d
\le20\,2^{d+15/2}d^{(25/4)d+29/4}.
\] For \(d\ge3\), \(20\,2^{15/2}<4096\le d^8\), while
\(2^d\le3^{d-1}\le d^{d-1}\): the first inequality starts at \(8\le9\)
and its ratio decreases by \(2/3\) with each increment of \(d\). Hence
\[
\widehat C_d^d\le d^{(29/4)d+57/4}\le d^{12d},
\] where the last inequality is exactly \(19d\ge57\). This proves (1).
\(\square\)

\paragraph{4.1. The direct geometric recovery
theorem}\label{the-direct-geometric-recovery-theorem}

The full proof in \url{geometry-rigidity.md} gives the following
stronger geometric input. For a centered symmetric full-dimensional body
\(Q\subset RB_2^d\), an orthogonal decomposition into \(m\) nonzero
blocks, its projection product \(P\), and
\(s=\mathbb E_{\nu_Q}\operatorname{dist}(X,\mathcal W)\), put \[
T=(2^{2d-1}dR s)^{1/(d-1)}.
\] Then \[
(1-mT)P\subset Q\subset P                                  \tag{9}
\] whenever \(mT<1\). In particular the Banach--Mazur excess to this
product is at most \(2mT\) when \(mT\le1/2\). If the factors have
dimensions one or two, the globally valid whole-class excess is at most
\(2mdT\). No inradius or inverse-basis norm occurs in these estimates.

For clarity, the mechanism and constants can be summarized
independently. For an orthogonal bipartition \(E\oplus F\), write
\(P_{E,F}\) for the two-factor projection product,
\(\lambda P_{E,F}\subset Q\) for its maximal centered homothetic
inclusion, and \(\tau=1-\lambda\). A support witness \(u,p\) has
\(u\cdot p=H=h_{P_{E,F}}(u)\) and \(h_Q(u)=(1-\tau)H\). Lift the two
block components of \(p\) to \(q_A,q_B\in Q\). Their differences from
\(p\) span a two-dimensional plane \(L\) with orthogonal axes in
\(E,F\). Its normal component \(u_L\) is nonzero.

Contract \(p,q_A,q_B\) by \(1-\eta\) toward any interior point of \(Q\),
with \(\eta=\tau/4\). The contracted witness remains outside \(Q\) by
support height at least \(\tau H/2\). In every resulting planar section,
the mixed-normal boundary length is therefore at least
\(\tau H/(2|u_L|)\). The planar assertion is proved by writing its
negative quadrant frontier as \(y=-g(x)\): convexity makes \(g'\)
nondecreasing, and \(\int\min(1,g')=\min(g(x)-x)\) is at least the
distance to the witness.

The transverse parameters fill a scaled copy of
\(\operatorname{proj}_{L^\perp}Q\) of measure
\(\eta^{d-2}|\operatorname{proj}_{L^\perp}Q|\). The coarea Jacobian on
the regular boundary is \(|\operatorname{proj}_L n|\), so \[
I_{E,F}:=\int_{\partial Q}
\min(|\operatorname{proj}_E n|,|\operatorname{proj}_F n|)\,d\mathcal H^{d-1}
\ge\eta^{d-2}|\operatorname{proj}_{L^\perp}Q|
      \frac{\tau H}{2|u_L|}.
\] Every planar fiber has area at most \(4RH/|u_L|\), by the same
support slab and the outer-radius bound. Fubini cancels the transverse
projection and support denominator, giving \[
\tau^{d-1}\le2^{2d-1}R I_{E,F}/|Q|.
\] For a single-block bipartition, its normal integrand is at most
\(\operatorname{dist}(n,\mathcal W)\). Homogeneity of distance cancels
the cone weight, so the latter integral is exactly \(d|Q|s\), and every
single-block homothetic loss is at most \(T\).

The single-block inclusions multiply: successively replace coordinates
by their desired projection-product values, scaled by the preceding
inclusion factors. The resulting point is in \(Q\) and has all
coordinates scaled by \(\prod_j\lambda_j\). Thus
\((\prod_j\lambda_j)P\subset Q\), with
\(\prod_j\lambda_j\ge1-\sum_j(1-\lambda_j)\ge1-mT\), proving (9). The
full note supplies the planar lemma and coarea regularity details for
arbitrary convex boundaries, including flat facets and nonatomic laws.

\paragraph{4.2. Complete assembly and numerical constant for Theorem
2}\label{complete-assembly-and-numerical-constant-for-theorem-2}

Retain \(M_d=\widehat L_d=5d^2a_d^{-2}\sqrt{B_d}\), and define \[
G_d=(2^{2d-1}d^2M_d)^{1/(d-1)},\qquad C'_d=2d^2G_d,
\qquad\beta_d=\frac1{3(d-1)}.
\] The coordinate body from (7) has \(R=d\), \(m\le d\), and
\(s\le M_d\delta^{1/3}\). Therefore \(T\le G_d\delta^{\beta_d}\). For \[
0<\delta\le\delta'_d:=(1/(2dG_d))^{3(d-1)},
\] we have \(mT\le1/2\). The matching condition \(q\le1\) also holds:
\(M_d\ge\sqrt{B_d}\), hence \(G_d^{d-1}\ge\sqrt{B_d}\), and \[
q\le B_d/(2dG_d)^{2(d-1)}\le1.
\] The small-deficit excess is consequently at most
\(2dG_d\delta^{\beta_d}\). At all larger deficits use the John bound
\(D-1\le d-1\); it is covered by the global coefficient since
\(C'_d(\delta'_d)^{\beta_d}=d\). At zero deficit, the exact matching
case and (9) give the product itself. This proves
\(D-1\le C'_d\delta^{\beta_d}\) in every case.

To bound \(C'_d\), the earlier estimate on \(B_d\) gives \[
M_d\le5\,2^{15/2}d^{(21/4)d+19/4},
\qquad
(C'_d)^{d-1}\le5\,2^{3d+11/2}d^{(29/4)d+19/4}.
\] For \(d\ge3\), \(8\le d^2\), and \(5\,2^{11/2}<243\le d^5\). Thus \[
(C'_d)^{d-1}\le d^{(37/4)d+39/4}\le d^{19(d-1)},
\] where the final inequality is exactly \(39d\ge115\), valid already at
\(d=3\). Hence \(C'_d\le d^{19}\), proving (8). \(\square\)

\subsubsection{5. What is genuinely sharp, and what remains
open}\label{what-is-genuinely-sharp-and-what-remains-open}

The accompanying complete notes isolate different hypotheses.

\begin{itemize}
\tightlist
\item
  \url{halfmass.md} proves that the cone-specific power \(2/3\), its
  near-half-mass tail power two, and the contact-slab area power two are
  sharp for actual cone laws in every fixed \(d\ge3\). The bodies are a
  doubly truncated three-cube times intervals, already in John position.
\item
  \url{matching.md} proves that the improved power \(1/2\) is sharp for
  bounded laws with determinant mass bounded below, even over every
  choice of basis and every line/plane decomposition. Its first
  obstruction is an arbitrary even law. The separate
  \url{global-next.md} strengthens the sharpness statement to actual
  cone laws and proves a geometric all-product distance obstruction for
  that family.
\item
  \url{geometry-next.md} proves that the cap exponent \(1/d\) is sharp
  under volume-only information. With the stronger off-block normal
  measurement, truncated cubes instead exclude powers above \(1/(d-1)\).
  The general recovery theorem with power \(1/(d-1)\), now proved
  directly in \url{geometry-rigidity.md}, closes this intermediate
  geometric gap. It is sharp for the stronger normal input.
\end{itemize}

These sharpness examples do not establish optimality of (8): independent
losses need not be simultaneously sharp on the same cone law. They rule
out improving the first two isolated lemmas under their present
hypotheses, and explain why a further gain should retain more geometry
or couple the stages.

The universally valid weighted relation \[
\mathbb E[\phi g]\le(d+1)A\delta,
\quad g=1/2-\max_i(|c_i|/\sum_j|c_j|),
\] with \(g=0\) on the singular event \(\sum_j|c_j|=0\), retains
information lost by the unweighted fourth-cofactor moment. On the sharp
half-mass family, the actual mixed-volume support surplus is linear in
\(\delta\). However, inserting an extra global containment factor into
the support-surplus estimate is false: the explicit near-coordinate cuts
in \url{geometry-next.md} give a counterexample. No unproved coupled
support-surplus theorem is used in either bound.

No quantitative inverse-Minkowski theorem is a proof input. The direct
section/coarea proof supplies the geometric estimate used here.

\subsubsection{6. Dependency and verification
scope}\label{dependency-and-verification-scope}

The complete baseline proof is included in
\url{dependencies/completed-proof.md}. Its corrected text has SHA-256
\nolinkurl{6a31c8d2dff76b3f73b862a0b6ef2ea192afc9de23b3605208b465ee3a5c1193};
the immutable original has SHA-256
\nolinkurl{08f6f55c6a694da05d348d3dcbda18be186044fd39570c51bf786507aaef862c}.
The baseline pins the v2--v4 sources of \texttt{mxym/math} to
\nolinkurl{6785c1c830f8e19e2eb07b0bb89f4d475a8b154a}. The exact dependency
map remains applicable. Bibliographic attributions and comparisons are
corrected separately; this package includes the corrected baseline
mathematical proof in \texttt{dependencies/completed-proof.md}. Neither
the lower-end simplex modulus nor the v5 spectral recursion is imported
here.

The new rational programs check the polynomial identities and exact
examples described in their adjacent notes. They use explicit exceptions
and retain checks under Python optimization. They do not prove the
general integral, surface-measure, or convex-geometric assertions. The
written proof and its independent model audit are separate from these
finite regressions. No human peer review, proof-assistant verification,
or novelty claim is made.

The optimal final exponent remains unresolved between \(1/[3(d-1)]\) and
\(1/d\). The isolated intermediate examples do not establish optimality
of the composed bound. The original completed release is preserved.

\newpage

\section{II Retained baseline proof and whole class
obstruction}\label{ii-retained-baseline-proof-and-whole-class-obstruction}

\subsubsection{1. The theorem and explicit
constants}\label{the-theorem-and-explicit-constants}

Let \(K\subset\mathbb R^d\) be any full-dimensional centrally symmetric
convex body, with its center translated to zero. There is no polytope,
smoothness, strict-convexity, absolute-continuity, or positive-curvature
assumption. Retain the entry005 invariant \[
 a(K)=\left(\frac d{d+1}\right)^d\frac{R(\mathcal P K)}{R(K)}-1,
 \qquad R(K)=\frac{|\Pi K|}{|K|^{d-1}}.
\] Let \(\mathcal E_d\) consist of all linear images of Cartesian
products of centrally symmetric factors of dimensions one and two. For
centered symmetric bodies use \[
 d_{\rm BM}(K,L)=\inf\{\lambda\ge1:\ TL\subseteq K\subseteq\lambda TL,
                  \ T\in\mathrm{GL}(d)\},
 \qquad D(K,\mathcal E_d)=\inf_{E\in\mathcal E_d}d_{\rm BM}(K,E).
\] Write \(\delta=\tfrac12-a(K)\), which is nonnegative by v3.

\textbf{Theorem A. Explicit power stability for the entire equality
class.} For \(d=1,2\), \(D(K,\mathcal E_d)=1\). For every \(d\ge3\),
define \[
 \begin{aligned}
 N&=d+1, &\kappa_d&=\frac{\pi^{d/2}}{\Gamma(1+d/2)}, &a_d&=d^{-2d},\\
 Z_d&=\frac{2^{d+9}N\sqrt d\binom d3}{\kappa_d},
 &F_d&=13N+3Z_d, & B_d&=\frac{4NF_d}{a_d},\\
 L_d&=6d^3a_d^{-4}B_d^{1/4},
 &H_d&=d(\sqrt d+1)L_d,
 &C_d&=2d(2dH_d)^{1/d}.
 \end{aligned}
\] Then, for every such \(K\), \[
 \boxed{D(K,\mathcal E_d)-1\le C_d\delta^{1/(6d)}.}                 \tag{1.1}
\] In particular, \(C_d\le d^{15}\), so the following simple numerical
version holds in every dimension \(d\ge3\): \[
 \boxed{D(K,\mathcal E_d)-1\le
        d^{15}\bigl(\tfrac12-a(K)\bigr)^{1/(6d)}.}              \tag{1.1a}
\] Equivalently, \(\delta\le(\varepsilon/d^{15})^{6d}\) guarantees a
Banach--Mazur factor at most \(1+\varepsilon\). All constants are
numerical functions of dimension. They are deliberately conservative. A
more informative small-deficit bound is \[
 D(K,\mathcal E_d)
 \le\frac1{1-(2dH_d\delta^{1/6})^{1/d}},                         \tag{1.2}
\] whenever \[
 0<\delta\le\delta_d^*:=
       \left(\frac{2^{-d}}{2dH_d}\right)^6.                       \tag{1.3}
\] At zero deficit the v4 equality classification applies, or the
zero-defect versions of the lemmas below give the same conclusion
directly.

\textbf{Theorem B. A dimensional obstruction.} For \(d\ge3\) and
\(0<t\le1\), set \[
 K_t=[-1,1]^d\cap\{x:\ |x_1+\cdots+x_d|\le d-t\},
 \qquad c_d=\frac12-2^{1-d}.
\] Then \[
 1+\frac{t}{5d^2}\le D(K_t,\mathcal E_d)\le\frac d{d-t},          \tag{1.4}
\] and \[
 \frac12-a(K_t)
 =c_d\frac{2^d-2t^{d-1}/(d-1)!}{d2^d-2t^d/(d-1)!}
       \frac{2t^d}{(d-1)!2^d+2(d-1)t^{d-1}}
 \sim\frac{c_d}{d!2^{d-1}}t^d.                                  \tag{1.5}
\] Consequently no universal fixed-dimensional power exponent greater
than \(1/d\) is possible. No positive exponent independent of dimension
can work, even when its constants are allowed to depend on dimension.
The optimal exponent remains between \(1/(6d)\) and \(1/d\); neither
theorem identifies it.

The proof combines three effective estimates: a cone-specific
contact-slab bound, a quantitative matching of basis coordinates, and a
mixed-volume containment argument. The hypothesis and equality class
differ from those of the lower-end simplex stability problem.

\subsubsection{2. Cone law, contact identity, and the exact upper-end
defect}\label{cone-law-contact-identity-and-the-exact-upper-end-defect}

Sections 2--5 assume \(d\ge3\); dimensions one and two were settled in
the statement.

The v3 cone law is the pushforward of \(h_K(u)dS_K(u)/(d|K|)\) by
\(u\mapsto u/h_K(u)\). Denote it by \(\nu_K\). For independent samples
define \[
 A=\mathbb E|\det(X_1,\ldots,X_d)|,
 \qquad
 B=\mathbb E\left|\det\begin{pmatrix}X_1&\cdots&X_N\\1&\cdots&1\end{pmatrix}\right|.
\] These determinant expectations are classical zonoid volumes. In the
normalization \(Z(\nu)=\mathbb E[0,X]\) and
\(\widehat Z(\nu)=\mathbb E[0,(1,X)]\), Koshevoy--Mosler, \emph{Lift
zonoids, random convex hulls and the variability of random vectors},
Bernoulli 4 (1998), 377--399, Section 4, gives \[
 |Z(\nu)|=A/d!,\qquad |\widehat Z(\nu)|=B/(d+1)!.
\] Consequently the cone-law identity used here reads \[
 a(K)=|\widehat Z(\nu_K)|/|Z(\nu_K)|.
\] The determinant representation is established background; the
restriction to cone laws and the endpoint inequality, equality analysis
and stability are separate assertions. See the
\href{https://wisostat.uni-koeln.de/sites/dep_econometrics/pdf_publikationen/1998_KoshevoyM98.pdf}{primary
author PDF}, printed p.~387, Theorem 4.1 and equation (4.3).

We import the v3 identity \(a(K)=B/(NA)\), affine invariance, and the
balanced Rademacher inequality. These hold for arbitrary convex bodies,
by the surface-area formulation and the continuity argument proved
there. The law is even for symmetric \(K\).

Put \(K\) in symmetric John position: \[
 B_2^d\subseteq K\subseteq\sqrt d B_2^d.                          \tag{2.1}
\] Then \(\nu_K\) is supported in the unit ball and \(|K|\ge\kappa_d\).
The law can also be sampled by a regular boundary contact
\(Y\in\partial K\), with density \[
 \frac{\langle Y,n_K(Y)\rangle}{d|K|}\,d\mathcal H^{d-1}(Y),
 \qquad X=\frac{n_K(Y)}{\langle Y,n_K(Y)\rangle}.
\] The nonregular boundary has surface measure zero. The divergence
theorem gives \[
 \langle Y,X\rangle=1,
 \qquad \mathbb E[Y X^{\mathsf T}]=I/d.                          \tag{2.2}
\] For every unit vector \(v\), since \(|Y|\le\sqrt d\), \[
 \mathbb E|\langle v,X\rangle|\ge\frac1{d\sqrt d}.
\] Conditioning on successive samples and using distance to their span,
this implies \[
 A=\mathbb E\|X_1\wedge\cdots\wedge X_d\|
 \ge d^{-3d/2}\ge a_d.                                         \tag{2.3}
\] Indeed, for every fixed proper subspace choose a unit vector
perpendicular to it; the conditional expected distance of the next
sample to that subspace is at least \(d^{-3/2}\). Multiplying these
lower bounds proves (2.3), with dependent partial tuples contributing
zero harmlessly. Hadamard gives \(A\le1\).

For an \(N\)-tuple, let \(c_i\) be its signed horizontal cofactors and
put \(M=\sum|c_i|\). Their relation is \(\sum c_iX_i=0\). Since the
points lie on \(\partial K^\circ\), the Minkowski norm with unit ball
\(K^\circ\) gives \(|c_i|\le M/2\). On \(M>0\), sort the normalized
magnitudes as \(t_1\ge\cdots\ge t_N\), and put \[
 \mathscr D(t)=\frac12-
        \mathbb E_\epsilon\left|\sum_i\epsilon_i t_i\right|,
 \qquad g=\frac12-t_1.
\] Define expressions involving \(t\) to be zero on \(M=0\). Sign
averaging the lifted determinant, using evenness, gives the exact
identity \[
 \Delta:=\mathbb E[M\mathscr D(t)]=NA\delta,
 \qquad \mathbb E M=NA\le N.                                   \tag{2.4}
\] This includes singular tuples. Let \(\phi=Mt_4\) be the
fourth-largest absolute cofactor, set to zero if all cofactors vanish.

\subsubsection{3. Robust balanced equality and the cone-specific
half-mass
gate}\label{robust-balanced-equality-and-the-cone-specific-half-mass-gate}

\textbf{Lemma 3.1.} For balanced normalized coefficients as above, \[
 \mathscr D(t)\ge\min\{t_4/2,g/4\}.                             \tag{3.1}
\] Their \(\ell_1\)-distance to the entire coefficient equality set (at
most three nonzero entries, or one entry equal to one half) is at most
\(8(N-3)\mathscr D(t)\).

\textbf{Proof.} On the capped simplex
\(\{s_i\ge0,\ \sum s_i=1,\ s_i\le1/2\}\), the convex function
\(\mathbb E_\epsilon|\sum\epsilon_i s_i|\) is at most one half by the v3
inequality. Let \(y=(1/4,1/4,1/4,1/4,0,\ldots)\), whose value is
\(3/8\). Set \(\alpha=\min(4t_4,2g)\). If positive, \(\alpha<1\), and
\((t-\alpha y)/(1-\alpha)\) remains in the capped simplex: nonnegativity
uses \(\alpha\le4t_4\); for the first four entries the upper cap is
equivalent to \(t_i+\alpha/4\le1/2\); for the others it follows from
\(\alpha\le2g\). Convexity yields a defect at least \(\alpha/8\),
proving (3.1). If \(\alpha=0\), (3.1) is immediate. To reach the
half-mass branch, increase \(t_1\) by \(g\) and subtract that mass from
the other entries, costing \(2g\). To reach three-point support, remove
the tail of mass \(q\le(N-3)t_4\) and redistribute it in the first three
entries without exceeding the cap, costing \(2q\). Their available
capacity is \(1/2+q\). Use the first construction when the minimum in
(3.1) is \(g/4\), and the second otherwise. This proves the distance
bound. \(\square\)

The half-mass branch cannot be dropped: \((1/2,1/6,1/6,1/6)\) has zero
defect and four positive entries. Exact exposedness in v4 removes it at
equality, but does not supply a quantitative margin. We instead use
contact geometry.

\textbf{Lemma 3.2. Corner-slab surface area.} If \(Q\subseteq[-1,1]^d\)
is any convex body, the regular boundary part where some three
coordinates are within \(r>0\) of independently chosen signed cube faces
has area at most \[
 d2^{d+1}\binom d3r^2.                                         \tag{3.2}
\]

\textbf{Proof.} Fix a triple and its signs and write \(\ell=\min(r,2)\).
Each coordinate projection of the relevant boundary part is contained in
a box of \((d-1)\)-volume at most \(2^{d-3}\ell^2\). If the deleted
coordinate is outside the triple, the sharper bound is
\(2^{d-4}\ell^3\), which is no larger. For any measurable boundary part
\(H\), the area formula and convexity give \[
 \int_H |n_i|\,d\mathcal H^{d-1}\le2|\pi_iH|.
\] Almost every fiber has at most an upper and a lower boundary point
with nonzero normal component. Parallel boundary segments have zero such
component. Since \(\sum_i|n_i|\ge1\), summing gives area at most
\(d2^{d-2}r^2\) for one signed triple. There are \(8\binom d3\) choices,
proving (3.2), also when \(r\ge2\). \(\square\)

\textbf{Lemma 3.3. Weighted exclusion of long near-half-mass circuits.}
For \(0<\eta\le1/4\) and \(\rho>0\), \[
 \mathbb E\bigl[M\mathbf1_{\{t_4>\rho,\ g\le\eta\}}\bigr]
 \le Z_d(\eta/\rho)^2.                                        \tag{3.3}
\] The event is understood on \(M>0\), as already implied by
\(t_4>\rho\).

\textbf{Proof.} Take a union over the distinguished index \(k\) with
\(|c_k|/M\ge1/2-\eta\), and condition on the other \(d\) samples,
arranged as columns of a matrix \(U\). If \(\det U=0\), this event at
\(k\) is impossible. Otherwise put \(D_0=|\det U|\), \(X_k=U\alpha\),
and \(S=\sum|\alpha_j|\). Cofactor identities give \[
 |c_k|=D_0,\quad |c_j|=D_0|\alpha_j|,
 \quad M=D_0(1+S),\quad \frac{|c_k|}M=\frac1{1+S}.
\] On this event \(M\le4D_0\) and \(S-1\le8\eta\). The condition
\(t_4>\rho\) supplies at least three other coordinates with
\(|\alpha_j|>\rho\).

Set \(Q=U^{\mathsf T}K\). Because the columns of \(U\) belong to
\(K^\circ\), \(Q\subseteq[-1,1]^d\) and \(|Q|=D_0|K|\). Under a primal
linear map \(T\), normalized cone mass is preserved and the scaled
normal becomes \(T^{-\mathsf T}X\); hence \(\nu_Q=(U^{-1})_*\nu_K\).
This follows directly from the regular-boundary formulas \[
 n'=\frac{T^{-\mathsf T}n}{|T^{-\mathsf T}n|},\quad
 dH'=|\det T|\,|T^{-\mathsf T}n|\,dH,\quad
 h'=h/|T^{-\mathsf T}n|.
\] For every regular contact \(z\in\partial Q\) generating the
conditional sample \(\alpha\), \(\langle\alpha,z\rangle=1\), so \[
 \sum_j|\alpha_j|\bigl(1-\operatorname{sgn}(\alpha_j)z_j\bigr)
 =S-1\le8\eta.
\] The three substantial coordinates therefore put \(z\) within
\(r=8\eta/\rho\) of three signed cube faces. Lemma 3.2 bounds their
surface area. The boundary cone density on \(Q\) is
\(\langle z,n_Q(z)\rangle/(d|Q|)\le\sqrt d/(dD_0|K|)\), so the
conditional probability is at most \[
 \frac{2^{d+1}\sqrt d\binom d3}{D_0|K|}(8\eta/\rho)^2.
\] Multiplication by \(M\le4D_0\) cancels the determinant exactly.
Integrate the conditioning variables, sum over \(N\) indices, and use
\(|K|\ge\kappa_d\). The resulting constant is \(Z_d\). \(\square\)

This cancellation treats arbitrarily ill-conditioned conditional
matrices. The estimate uses regular primal contacts, not strong
exposedness of polar points, so atoms caused by flat faces and nonatomic
cone laws are both covered.

\textbf{Lemma 3.4. Effective short-cofactor estimate.} \[
 \boxed{\mathbb E\phi\le F_d\delta^{2/3}.}                       \tag{3.4}
\]

\textbf{Proof.} First suppose \(0<\delta<1/4\). For
\(\rho\in[\delta,1/4]\), put \(\eta=\delta^{1/3}\rho^{2/3}\le\rho\). In
the weighted tail \(T(\rho)=\mathbb E[M\mathbf1_{\{t_4>\rho\}}]\), the
part with \(g>\eta\) costs at most \(4\Delta/\eta\) by Lemma 3.1. The
other part costs at most \(Z_d(\eta/\rho)^2\) by Lemma 3.3. Using (2.4),
\[
 T(\rho)\le(4N+Z_d)\delta^{2/3}\rho^{-2/3}.
\] Since \(t_4\le1/4\), Tonelli gives
\(\mathbb E\phi=\int_0^{1/4}T(\rho)\,d\rho\). The interval below
\(\delta\) costs at most \(N\delta\le N\delta^{2/3}\); the remaining
integral of \(\rho^{-2/3}\) is at most 3. This gives (3.4). If
\(\delta\ge1/4\), use \(\mathbb E\phi\le N/4\le N\delta^{2/3}\). If
\(\delta=0\), for fixed \(\rho>0\) let \(\eta\downarrow0\) in the same
split, with \(\eta\le\min(1/4,2\rho)\), to get \(T(\rho)=0\). Tonelli
proves (3.4) at zero without compactness. \(\square\)

\subsubsection{4. A quantitative matching of line and plane
blocks}\label{a-quantitative-matching-of-line-and-plane-blocks}

The exact rank-two decomposition underlying this lemma has a matroid
antecedent: van der Pol--Walsh--Wigal,
\href{https://arxiv.org/html/2502.03673v2\#S5}{arXiv:2502.03673v2},
Lemma 5.2, classify \(U_{3,4}\)-minor-free matroids as direct sums of
rank-at-most-two matroids. The estimate below concerns small expected
cofactor magnitudes and supplies a conditioned basis and quantitative
approximation to such blocks.

\textbf{Lemma 4.1.} Let \(\nu\) be any Borel probability law supported
in the unit ball, with \(A\ge a\), \(0<a\le1\), and suppose its expected
fourth absolute cofactor is at most \(\varepsilon\). Put
\(q=4(d+1)\varepsilon/a\). If \(0<q\le1\), there are a basis matrix
\(U\) and a partition of its coordinate indices into singletons and
pairs, with coordinate block union \(\mathcal W\), such that \[
 |\det U|\ge a/2,\quad \|U\|\le\sqrt d,\quad
 \|U^{-1}\|\le2\sqrt d/a,
\] \[
 \mathbb E\operatorname{dist}(U^{-1}X,\mathcal W)
 \le6d^3a^{-4}q^{1/4}.                                       \tag{4.1}
\] For \(\varepsilon=0\), the same conclusion holds with zero distance.
All statements include nonatomic laws and singular sampled tuples.

\textbf{Proof.} For a deterministic candidate basis
\(b=(b_1,\ldots,b_d)\), define its section cost \[
 Q(b)=\int\phi(b,x)d\nu(x)
 +\sum_{j=1}^d\iint\phi(b_1,\ldots,\widehat b_j,\ldots,b_d,x,y)
                         d\nu(x)d\nu(y).
\] Tonelli gives \(\mathbb E_bQ(b)\le(d+1)\varepsilon\). Since
\(|\det b|\le1\), the event \(|\det b|\ge a/2\) has probability at least
\(a/2\). Markov gives \(\mathbb P(Q>q)\le a/4\). Thus choose \(U=(b_i)\)
with \(D:=|\det U|\ge a/2\) and \(Q(U)\le q\). The norm bounds follow
from unit-length columns and inverse rows, each of length at most
\(1/D\). Cramer's rule also gives \(|\alpha_i|\le1/D\) for
\(\alpha=U^{-1}X\).

Let \(s_3\) be the third-largest coordinate magnitude. The absolute
cofactors of \((b_1,\ldots,b_d,X)\) are \(D\) times
\((1,|\alpha_1|,\ldots,|\alpha_d|)\). If \(s_3\le1\), its fourth largest
is at least \(Ds_3\); otherwise it is at least \(D\), and \(s_3\le1/D\).
In both cases \[
 \mathbb E s_3\le q/D^2.                                      \tag{4.2}
\] Fix \(0<t\le1\) and set \(r=Dt^2/2\). Break coordinate ties by index.
For each pair \(ij\), let \(E_{ij}\) be the event that its coordinates
are the largest two, both have magnitude at least \(t\), and
\(s_3\le r\). Write \(p_{ij}=\nu(E_{ij})\).

For distinct \(i,j,k\), take \(x\in E_{ij}\), \(y\in E_{jk}\), with
coordinates \(\alpha,\beta\). In the tuple consisting of the basis
except \(b_j\), followed by \(x,y\), the cofactors attached to
\(x,y,b_i,b_k\) have, up to signs, the values \[
 D\beta_j,\quad D\alpha_j,\quad
 D(\alpha_i\beta_j-\alpha_j\beta_i),\quad
 D(\alpha_k\beta_j-\alpha_j\beta_k).
\] These are determinant identities without a rank assumption. The first
two have magnitude at least \(Dt\). Each of the last two has magnitude
at least \(D(t^2-r/D)=Dt^2/2\), since the unwanted coordinate is at most
\(r\) and the other factor at most \(1/D\). Hence \[
 p_{ij}p_{jk}\le\frac{2q}{Dt^2}.
\] Pairs with mass exceeding \(w=\sqrt{2q/D}/t\) form a matching. Take
its pairs as plane blocks and every remaining index as a line block.
Discarded pairs have total mass at most \(\binom d2w\).

Outside the discarded-pair events and \(\{s_3>r\}\), a point is within
\(\sqrt d\,t\) of its block: if its second coordinate is below \(t\),
use the block containing its largest coordinate; otherwise use its
matched pair, whose remaining coordinates are at most \(r\le t\). Every
point has coordinate norm at most \(\sqrt d/D\). By (4.2) and Markov, \[
 \mathbb E\operatorname{dist}(\alpha,\mathcal W)
 \le\sqrt d\,t+\frac{\sqrt d}{D}
 \left[\binom d2\frac{\sqrt{2q/D}}t+\frac{2q}{D^3t^2}\right].     \tag{4.3}
\] Take \(t=q^{1/4}\). Since \(q\le1\) and \(D\ge a/2\), the right side
is at most \[
 \sqrt d\,[1+4\binom d2a^{-3/2}+32a^{-4}]q^{1/4}
 \le6d^3a^{-4}q^{1/4}
\] for \(d\ge3\). This proves (4.1).

At \(\varepsilon=0\), choose \(Q(U)=0\) on the determinant-good event.
Equation (4.2) makes every sample two-sparse almost surely. If two
positive-mass two-coordinate classes shared an index, the displayed
omitted-basis tuple would have four nonzero cofactors on their product,
contradicting its zero section cost. Thus the positive-mass pairs form a
matching and the coordinate distance is zero. This also proves the zero
case. \(\square\)

The basis is selected with simultaneous control of all required
sections. Its condition number is bounded explicitly. No limiting choice
of a basis or unique product decomposition is assumed.

\subsubsection{5. From approximate normals to an actual product
body}\label{from-approximate-normals-to-an-actual-product-body}

\textbf{Lemma 5.1.} Suppose \(Q\) is a full-dimensional centered
symmetric convex body, \(Q\subseteq R B_2^d\), and
\(\mathbb R^d=\bigoplus_jW_j\) is an orthogonal sum. Set
\(\mathcal W=\bigcup_jW_j\) and \(P=\prod_j\operatorname{proj}_{W_j}Q\),
with \(m\) blocks. If
\(s=\mathbb E_{\nu_Q}\operatorname{dist}(X,\mathcal W)\), then \[
 |P|/|Q|\le(1+R(\sqrt m+1)s)^d.                               \tag{5.1}
\] Putting \(v=R(\sqrt m+1)s\), if \(2dv<1\), \[
 (1-(2dv)^{1/d})P\subseteq Q\subseteq P.                       \tag{5.2}
\]

The conversion from relative volume loss to containment is the standard
nested-body cap argument; compare Böröczky--Henk,
\href{https://www.renyi.hu/~carlos/cone-volume-stability.pdf}{\emph{Cone-volume
measure and stability}}, Lemma 3.1(i), printed p.~8. Central symmetry
permits the constant 2 in place of the centered-body constant e in that
result.

\textbf{Proof.} The projection product contains \(Q\), is symmetric, and
lies in \(\sqrt m R B_2^d\). Its support function agrees with \(h_Q\) on
every block. For \(x\in\partial Q^\circ\), choose a closest block
projection \(z\in\mathcal W\). Support functions are globally Lipschitz
with their circumradii, so \[
 0\le h_P(x)-1=h_P(x)-h_Q(x)
 \le R(\sqrt m+1)|x-z|.
\] The cone-law definition, homogeneity, and the first mixed-volume
formula give \[
 \frac{V(Q[d-1],P)}{|Q|}-1
 =\mathbb E_{\nu_Q}[h_P(X)-1]\le v.
\] Minkowski's first inequality gives (5.1). Consequently
\(\theta=1-|Q|/|P|\le1-(1+v)^{-d}\le dv\).

For completeness, a symmetric nested-body cap converts this relative
volume loss to containment. Let \(\lambda\) be maximal with
\(\lambda P\subseteq Q\); then \(0<\lambda\le1\). If \(\lambda<1\), a
touching supporting direction \(u\) and point \(p\in P\) satisfy
\(h_Q(u)=\lambda h_P(u)\), \(\langle u,p\rangle=h_P(u)\). The image \[
 \lambda p+(1-\lambda)\bigl(P\cap\{z:\langle u,z\rangle\ge0\}\bigr)
\] lies in \(P\) by convexity and outside \(\operatorname{int}Q\) by
support. Symmetry makes its volume \((1-\lambda)^d|P|/2\). Hence
\(\theta\ge(1-\lambda)^d/2\), proving (5.2). No inverse cone-volume
uniqueness theorem is used. \(\square\)

\textbf{Proof of Theorem A.} Use (2.1). Lemma 3.4 gives
\(\varepsilon=F_d\delta^{2/3}\). By (2.3), apply Lemma 4.1 with
\(a=a_d\), so \(q=B_d\delta^{2/3}\). For \(\delta\le\delta_d^*\),
\(q\le1\): indeed \(H_d\ge B_d^{1/4}\), so
\(B_d(2^{-d}/(2dH_d))^4\le1\). The resulting coordinate law satisfies \[
 \mathbb E\operatorname{dist}(U^{-1}X,\mathcal W)
 \le L_d\delta^{1/6}.
\] Apply the primal map \(U^{\mathsf T}\). Its body \(Q=U^{\mathsf T}K\)
lies in \(dB_2^d\), because \(\|U\|\le\sqrt d\), and its cone law is the
coordinate law just obtained. Its blocks are orthogonal coordinate lines
and planes. Lemma 5.1 gives \(v\le H_d\delta^{1/6}\), with a product
\(P\in\mathcal E_d\), proving (1.2). In the range (1.3), the quantity
\(r=(2dH_d\delta^{1/6})^{1/d}\) is at most \(1/2\), so \[
 D(K,\mathcal E_d)-1\le\frac r{1-r}
 \le2(2dH_d)^{1/d}\delta^{1/(6d)}.
\] For larger deficits, a universal elementary bound is
\(D(K,\mathcal E_d)\le d\). In John position,
\(d^{-1/2}[-1,1]^d\subseteq K\subseteq d^{1/2}[-1,1]^d\), and the cube
belongs to \(\mathcal E_d\). If \(\delta\ge\delta_d^*\), the constant in
(1.1) satisfies \[
 C_d\delta^{1/(6d)}\ge
 2d(2dH_d)^{1/d}(\delta_d^*)^{1/(6d)}=d,
\] which covers \(D-1\le d-1\). At zero use the zero cases of the lemmas
or v4. For \(d=1,2\), every symmetric body is itself an allowed factor.
This completes the proof. \(\square\)

For the stated simpler constant, the cube inscribed in the unit ball
gives \(\kappa_d\ge(2/\sqrt d)^d\). Using \(N\le2d\) and
\(\binom d3\le d^3\), one obtains \[
 Z_d\le2^{10}d^{d/2+9/2},\quad
 F_d\le2^{12}d^{d/2+9/2},\quad
 B_d\le2^{15}d^{(5/2)d+11/2},
\] \[
 H_d\le12\,2^{15/4}d^{(69/8)d+47/8},\qquad
 C_d^d\le2^d(24\,2^{15/4})d^{(77/8)d+55/8}.
\] Since \(2^d\le d^d\) and \(24\,2^{15/4}<512\le d^6\) for \(d\ge3\),
\[
 C_d^d\le d^{(85/8)d+103/8}\le d^{15d},
\] where the last inequality is \(35d\ge103\). This proves (1.1a).

\subsubsection{6. Proof of the dimensional
obstruction}\label{proof-of-the-dimensional-obstruction}

We give the full-class argument, since comparing only to the cube would
not establish a distance lower bound to arbitrary planar-factor
products.

Let \(D\in\mathcal E_d\) and suppose
\(K_t\subseteq D\subseteq(1+s)K_t\). Any defining Banach--Mazur sandwich
can be rescaled into this form, preserving \(\mathcal E_d\). Write \[
 D=F_1+\cdots+F_m,\qquad
 \mathbb R^d=W_1\oplus\cdots\oplus W_m,
 \qquad F_j\subset W_j,\quad \dim W_j\in\{1,2\}.
\] The factors are centered symmetric and full-dimensional in their
subspaces. Let \(P_j\) be the projection onto \(W_j\) along the
remaining summands, and set \(a_{ji}=h_{F_j}(e_i)\ge0\). Then \[
 P_j^2=P_j,\quad\operatorname{tr}P_j=\dim W_j,\quad\sum_jP_j=I,
 \qquad 1\le\sum_ja_{ji}\le1+s.                               \tag{6.1}
\] Every cube vertex \(\sigma\in\{\pm1\}^d\) except the two all-equal
ones belongs to \(K_t\). Since \(P_j\sigma\in F_j\), the nonnegative
quantities \(a_{ji}-\sigma_i(P_j\sigma)_i\) sum to at most \(s\).
Therefore \[
 |(P_j\sigma)_i-a_{ji}\sigma_i|\le s.                          \tag{6.2}
\] For any coordinate \(k\), assign both signs among the other
coordinates; this is possible since \(d\ge3\). Both vertices obtained by
changing only \(\sigma_k\) are retained. Subtract (6.2) for them to
obtain \[
 |(P_j-\operatorname{diag}(a_{j1},\ldots,a_{jd}))_{ik}|\le s.     \tag{6.3}
\] This controls all entries even when the original decomposition is
badly conditioned.

Suppose \(s\le1/(4d)\), and write \(p_{ji}=(P_j)_{ii}\). Idempotence
gives \(|p_{ji}(p_{ji}-1)|\le(d-1)s^2\). Thus some \(z_{ji}\in\{0,1\}\)
satisfies \[
 |p_{ji}-z_{ji}|\le2(d-1)s^2.
\] Inside \([0,1]\), this follows from \(q(1-q)\ge q/2\) for distance
\(q\le1/2\) to an endpoint; outside the interval it follows directly.
The integer trace and sum identities in (6.1), with errors less than one
half, give \(\sum_i z_{ji}=\dim W_j\) and \(\sum_jz_{ji}=1\). Hence
\(J_j=\{i:z_{ji}=1\}\) partitions the indices into sets of size one or
two. For \(i\notin J_j\), (6.3) gives \(a_{ji}\le s+2(d-1)s^2\).

Let \(u=(1,\ldots,1)\). Choose a retained vertex with plus signs on
\(J_j\) and a minus sign outside it. Equation (6.2) implies \[
 h_{F_j}(u)\ge\sum_i a_{ji}-3ds-4d(d-1)s^2
            \ge\sum_i a_{ji}-4ds.
\] Summing gives \(h_D(u)\ge d-4d^2s\), while containment in
\((1+s)K_t\) gives \(h_D(u)\le d-t+ds\). Therefore \(t\le5d^2s\). If
\(s>1/(4d)\), the same lower bound follows from \(t\le1\). This proves
the lower part of (1.4) for every product decomposition. The upper part
follows from \((1-t/d)[-1,1]^d\subseteq K_t\subseteq[-1,1]^d\).

For the exact deficit, the two removed corner simplices each have volume
\(t^d/d!\), and every coordinate facet has area
\(\alpha=2^{d-1}-t^{d-1}/(d-1)!\). The two new area-normal vectors are
\(\pm t^{d-1}u/(d-1)!\), with area times support number
\(\beta=(d-t)t^{d-1}/(d-1)!\). Put \[
 V_t=2^d-2t^d/d!,\quad D_t=d2^d-2t^d/(d-1)!,\quad
 p=\frac{2^d-2t^{d-1}/(d-1)!}{D_t},\quad
 r=\frac{2(d-t)t^{d-1}}{(d-1)!D_t}.
\] The cone law has direction pairs \(\pm e_i\), each with total
probability \(p\), and \(\pm u/(d-t)\), with total probability \(r\).
Within pairs signs are uniform; \(dp+r=1\). Direct type enumeration
gives \[
 A=d!p^{d-1}\left(p+\frac{dr}{d-t}\right).
\] Only \(N\)-tuples containing all \(d+1\) different direction types
have a nonzero Rademacher defect. Tuples with fewer than \(d\) types
have zero cofactors; those with exactly \(d\) types have a repeated
parallel direction and a two-term cofactor relation. For the
all-distinct tuple, cofactor magnitudes are 1 and \(d\) copies of
\(q=1/(d-t)\). Averaging the sign multiplying 1 gives
\(\mathbb E\max(1,|q\sum_{i=1}^d\epsilon_i|)\). For \(t\le1\), only the
two all-equal sign patterns exceed 1, so its value is
\(1+2^{1-d}(dq-1)\). The defect is \(c_dt/(d-t)\). The probability of
all-distinct types is \(N!p^dr\), giving \[
 \delta_t=\frac{c_dprt}{(d-t)p+dr}.
\] Substitute \(p,r\) to get (1.5). Also \(p\le1/d\), hence
\(\delta_t\le c_dt^d/(d!2^{d-1})\). Combining this with (1.4) excludes
every exponent greater than \(1/d\) as \(t\downarrow0\). In dimension
three, \[
 \delta_t=\frac{t^3(8-t^2)}{8(24-t^3)(4+t^2)}.
\] This proves Theorem B. \(\square\)

\subsubsection{7. Comparison with primary cone-volume and stability
literature}\label{comparison-with-primary-cone-volume-and-stability-literature}

The class \(\mathcal E_d\) is the familiar class of affine Cartesian
products of centrally symmetric intervals and planar bodies, called Weil
bodies in Saroglou, \href{https://arxiv.org/pdf/1409.4347v2}{\emph{On
the shape of a convex body with respect to its second projection body}},
pp.~2--4. That paper relates the class to Weil's polytopal
classification for homothety under the second projection-body operator
and gives a three-dimensional zonoid endpoint with symmetric-cylinder
equality (Theorem 1.2). These use a different functional from \(a(K)\);
the common extremal class alone does not make their statements
equivalent to entry005 v4's equality theorem or Theorem A here.

The relevant antecedent is Böröczky--Lutwak--Yang--Zhang,
\href{https://www.renyi.hu/~carlos/minkowski0-prob.pdf}{The logarithmic
Minkowski problem}, JAMS 26 (2013), 831--852. Its subspace concentration
condition and complementary-subspace clause characterize existence of
symmetric cone-volume measures. They do not supply the
near-short-circuit estimate (3.4), a quantitative matching into
dimensions at most two, or the present invariant's endpoint.

Böröczky--Henk,
\href{https://www.renyi.hu/~carlos/cone-volume-stability.pdf}{Cone-volume
measure and stability}, Adv. Math. 306 (2017), 24--50, Theorem 1.1 gives
complementary-factor stability with containment exponent \(1/(5d)\) from
almost saturated cone mass on an \textbf{exact} proper subspace; its
constants and smallness threshold are stated as dimension-dependent,
rather than numerical formulas. Its U-functional stability targets
parallelotopes. Our equality class includes all symmetric planar
factors, and our hypothesis is determinant magnitude, not the
U-functional's independence indicator. Thus that theorem is not an
equivalent endpoint result.

The exact-subspace premise cannot be inferred here: for even integers
\(p\to\infty\), smooth strictly convex \(\ell_p\) balls converge to the
cube and have \(a\to1/2\), yet their cone-volume measures give zero mass
to every proper linear subspace. On each open coordinate octant the
Gauss map is a diffeomorphism with positive curvature; the omitted
coordinate-zero parts have surface measure zero, proving absolute
continuity. Our method instead produces an expected distance to blocks
and recovers an actual product by Lemma 5.1.

Böröczky--De, \href{https://arxiv.org/abs/2101.03395}{Stable solution of
the logarithmic Minkowski problem in the case of hyperplane symmetries},
J. Differential Equations 298 (2021), 298--322, obtains quantitative
inverse cone-volume estimates under Coxeter symmetry and a tube
concentration margin. The reducible-group margin degenerates at
concentration on a proper group-invariant product block, and the
symmetry hypothesis does not cover arbitrary symmetric planar factors.
Their anisotropic-squeeze example also explains why normalization is
necessary. Abdallah--Mérigot, \href{https://arxiv.org/abs/1402.5010}{On
the reconstruction of convex sets from random normal measurements},
gives a surface-area inverse estimate of order \(1/d\) with area and
rotundity control. Cone-volume data additionally contains the unknown
support function; Lemma 5.1 avoids solving that inverse problem.

Böröczky--De also discuss volume-normalized unconditional bodies
obtained by truncating vertices of a cube
(\href{https://real.mtak.hu/196592/1/1-s2.0-S0022039621004393-main.pdf}{published
version}, p.~302). This provides related geometric context. Theorem B
independently computes the upper-end scalar deficit \(1/2-a(K_t)\) and
proves a distance lower bound against every affine product of symmetric
one- and two-dimensional bodies. A comparison with one fixed cube would
not establish that full-class lower bound.

Saroglou's \href{https://doi.org/10.1112/S0025579311001860}{Volumes of
projection bodies of some classes of convex bodies} is a primary
projection-body comparison. Only its publication record and abstract
were available; a full-text comparison remains incomplete. The recent
Liu--Xiong--Yang mixed-volume stability paper, DOI
\href{https://doi.org/10.1016/j.aim.2026.110886}{10.1016/j.aim.2026.110886},
is relevant to broader mixed-volume stability questions. No theorem from
either of these two papers is imported as a dependency. The present
literature comparison is confined to the primary statements described
above.

\subsubsection{8. Dependencies, reproducibility, and open
questions}\label{dependencies-reproducibility-and-open-questions}

The affine calculus and cone-law representation are inherited from entry
005 versions 2--3; the exact equality class is the version 4 theorem.
The pinned sources are listed below. The quantitative contact-slab,
matching, and product recovery arguments in Sections 3--5 are supplied
in full here. In particular, the lower-end simplex modulus and the
balanced-recursion spectral theorem are separate results and are not
inputs to this proof.

The accompanying exact programs test 4,930 balanced rational coefficient
multisets, 750 omitted-basis cofactor identities, and 16
corner-truncation facet-minor calculations. They use only integers and
rational arithmetic, and explicit exceptions keep the checks active
under Python optimization. These finite regressions supplement the
written analytic proof. The package records source hashes and separates
regenerated outputs from shipped reference reports.

The constants are deliberately coarse. The optimal universal exponent
remains between \(1/(6d)\) and \(1/d\) for \(d\ge3\). Sharpening the
matching estimate and the conversion of normal concentration to product
containment may reduce the gap.

\subsubsection{References and
provenance}\label{references-and-provenance}

\begin{itemize}
\tightlist
\item
  {[}005v2{]} mxym, \emph{Projection-volume calculus for joins and
  Cartesian products}, entry 005, version 2. Definitions, affine
  invariance, and product calculus.
  \href{https://github.com/mxym/math/blob/6785c1c830f8e19e2eb07b0bb89f4d475a8b154a/preprints/005-simplex-product-optimum/v2/paper.md}{Pinned
  source}.
\item
  {[}005v3{]} mxym, \emph{Random-determinant rigidity, sharp symmetric
  cone bounds, and spectral nonattainment}, entry 005, version 3.
  Cone-law representation and symmetric inequality.
  \href{https://github.com/mxym/math/blob/6785c1c830f8e19e2eb07b0bb89f4d475a8b154a/preprints/005-simplex-product-optimum/v3/paper.md}{Pinned
  source}.
\item
  {[}005v4{]} mxym, \emph{Equality in the symmetric projection-cone
  bound}, entry 005, version 4. The equality class and its qualitative
  stability.
  \href{https://github.com/mxym/math/blob/6785c1c830f8e19e2eb07b0bb89f4d475a8b154a/preprints/005-simplex-product-optimum/v4/paper.md}{Pinned
  source}.
\item
  {[}BLYZ{]} K.J. Böröczky, E. Lutwak, D. Yang, and G. Zhang, \emph{The
  logarithmic Minkowski problem}, Journal of the American Mathematical
  Society 26 (2013), 831--852.
  \href{https://doi.org/10.1090/S0894-0347-2012-00741-3}{DOI}.
\item
  {[}BH{]} K.J. Böröczky and M. Henk, \emph{Cone-volume measure and
  stability}, Advances in Mathematics 306 (2017), 24--50.
  \href{https://doi.org/10.1016/j.aim.2016.10.005}{DOI}.
\item
  {[}BD{]} K.J. Böröczky and A. De, \emph{Stable solution of the
  logarithmic Minkowski problem in the case of hyperplane symmetries},
  Journal of Differential Equations 298 (2021), 298--322.
  \href{https://doi.org/10.1016/j.jde.2021.07.002}{DOI}.
  \href{https://arxiv.org/html/2101.03395v2}{Primary preprint, version
  2}.
\item
  {[}AM{]} H. Abdallah and Q. Mérigot, \emph{On the reconstruction of
  convex sets from random normal measurements}.
  \href{https://arxiv.org/abs/1402.5010}{Primary preprint}.
\item
  {[}S{]} C. Saroglou, \emph{Volumes of projection bodies of some
  classes of convex bodies}, Mathematika 57 (2011), 329--353.
  \href{https://doi.org/10.1112/S0025579311001860}{DOI}.
\item
  {[}LXY{]} Y.-D. Liu, G. Xiong, and K.-W. Yang, \emph{Sharp
  quantitative stability for the Minkowski first inequality via a
  quadratic estimate for the cone-volume measure}, Advances in
  Mathematics 492 (2026), 110886.
  \href{https://doi.org/10.1016/j.aim.2026.110886}{DOI}.
\item
  {[}KM{]} G. Koshevoy and K. Mosler, \emph{Lift zonoids, random convex
  hulls and the variability of random vectors}, Bernoulli 4 (1998),
  377--399.
  \href{https://wisostat.uni-koeln.de/sites/dep_econometrics/pdf_publikationen/1998_KoshevoyM98.pdf}{Primary
  author PDF}, Section 4, Theorem 4.1 and equation (4.3).
\item
  {[}VPWW{]} J. van der Pol, Z. Walsh, and M.C. Wigal, \emph{Turán
  densities for matroid basis hypergraphs}, arXiv:2502.03673v2, 10 March
  2025. \href{https://arxiv.org/html/2502.03673v2\#S5}{Primary
  manuscript}, Section 5, Lemma 5.2 (printed p.~15).
\item
  {[}SS{]} C. Saroglou, \emph{On the shape of a convex body with respect
  to its second projection body}, arXiv:1409.4347v2.
  \href{https://arxiv.org/pdf/1409.4347v2}{Primary manuscript}, pp.~2--4
  and Theorem 1.2. It cites W. Weil, \emph{Über die Projektionenkörper
  konvexer Polytope}, Archiv der Mathematik 22 (1971), 664--672; that
  original article was not obtained for this comparison.
\item
  {[}Sch{]} R. Schneider, \emph{Convex Bodies: The Brunn--Minkowski
  Theory}, second expanded edition, Cambridge University Press, 2014.
  Support functions, surface area, John normalization, and mixed
  volumes.
\end{itemize}

Versions2--4 credit the projection-volume product and simplex
computations inherited from OpenAI family 088, pinned at commit
\nolinkurl{adc7f1241b42e322a6451854ab7e4b4c146bf78a}. This supplement uses
the inherited identities through the explicitly cited entry 005 sources.

\newpage

\section{III Weighted rank two
matching}\label{iii-weighted-rank-two-matching}

\subsection{A sharp square-root matching theorem for bounded
laws}\label{a-sharp-square-root-matching-theorem-for-bounded-laws}

This note strengthens the matching lemma in the completed upper-end
stability supplement. It is separate research; it does not alter that
supplement.

\subsubsection{Statement and
conventions}\label{statement-and-conventions}

Let \(d\ge3\), let \(\nu\) be a Borel probability measure supported in
the Euclidean unit ball, and let \(X_1,\ldots,X_{d+1}\) be independent
samples. Write \[
 A=\mathbb E|\det(X_1,\ldots,X_d)|.
\] For a tuple of \(d+1\) vectors, let \(\phi\) denote the
fourth-largest of the \(d+1\) absolute \(d\)-dimensional cofactors.
Suppose \[
 A\ge a>0,\qquad \mathbb E\phi\le\varepsilon,
 \qquad q=4(d+1)\varepsilon/a.
\] Hadamard's inequality implies \(a\le1\).

\textbf{Theorem 1 (square-root matching).} There are a basis matrix
\(U\) whose columns lie in the support of \(\nu\), and a partition of
its coordinates into singletons and pairs, such that the union
\(\mathcal W\) of the resulting coordinate blocks satisfies \[
 |\det U|\ge a/2,\qquad \|U\|\le\sqrt d,
 \qquad \|U^{-1}\|\le2\sqrt d/a,
\] and, if \(0\le q\le1\), \[
 \mathbb E\operatorname{dist}(U^{-1}X,\mathcal W)
 \le5d^2a^{-2}\sqrt q.                                      \tag{1}
\] The conclusion includes nonatomic laws and singular sampled tuples.
The power \(1/2\) cannot be increased under these hypotheses, even for
even finite laws and fixed dimension.

\subsubsection{Proof of Theorem 1}\label{proof-of-theorem-1}

For a candidate sampled basis \(b=(b_1,\ldots,b_d)\), put \[
 Q(b)=\int\phi(b_1,\ldots,b_d,x)\,d\nu(x)
 +\sum_{j=1}^d\iint
 \phi(b_1,\ldots,\widehat b_j,\ldots,b_d,x,y)
 \,d\nu(x)\,d\nu(y).
\] Tonelli's theorem gives \(\mathbb E_bQ(b)\le(d+1)\varepsilon\). Since
\(D(b)=|\det b|\le1\) and its expectation is at least \(a\), the event
\(D(b)\ge a/2\) has probability at least \(a/2\). If \(\varepsilon>0\),
Markov's inequality gives \(\mathbb P(Q(b)>q)\le a/4\). These two events
therefore intersect. If \(\varepsilon=0\), \(Q(b)=0\) almost everywhere
and the determinant-good event again intersects it. Fix such a basis
\(U=(b_i)\), and write \[
 D=|\det U|\ge a/2,\qquad R=D^{-1}.
\] Each inverse row has length at most \(D^{-1}\), by its cofactor-wedge
formula and the unit-length bound on the other columns. Consequently
\(\|U^{-1}\|\le\sqrt d/D\), \(\|U\|\le\sqrt d\), and the coordinates
\(\alpha=U^{-1}X\) obey \(\max_i|\alpha_i|\le R\).

Let \(s_3(\alpha)\) be the third-largest coordinate magnitude. The
absolute cofactors of \((b_1,\ldots,b_d,X)\) are \(D\) times
\((1,|\alpha_1|,\ldots,|\alpha_d|)\). If \(s_3\le1\), the fourth largest
is at least \(Ds_3\); otherwise it is at least \(D\), and
\(s_3\le D^{-1}\). Thus \[
 \mathbb Es_3\le Q_0/D^2\le q/D^2,                           \tag{2}
\] where \(Q_0\) is the one-sample term of \(Q(U)\).

For every \(\alpha\), retain its two largest coordinates, breaking ties
by index, and set all others to zero. Denote the resulting vector by
\(\widetilde\alpha\). This is a Borel map, it does not enlarge the
coordinate bounds, and \[
 \|\alpha-\widetilde\alpha\|_\infty\le s_3(\alpha),
 \qquad |\alpha-\widetilde\alpha|\le\sqrt d\,s_3(\alpha).    \tag{3}
\]

We first control the two-sample sections after this truncation. For
independent coordinate vectors \(\alpha,\beta\), the absolute cofactors
of the tuple omitting \(b_j\) and appending \(U\alpha,U\beta\) are the
absolute values of \[
 D\alpha_j,\quad D\beta_j,\quad
 D(\alpha_i\beta_j-\alpha_j\beta_i)\quad(i\ne j),             \tag{4}
\] up to signs and reordering. This is the polynomial determinant
identity obtained by applying \(U^{-1}\) to every column.

Replacing both vectors by their truncations changes either linear entry
by at most \(D(s_3(\alpha)+s_3(\beta))\). For every bilinear entry,
telescoping the two products and using all coordinate bounds
\(R=D^{-1}\) gives a change at most \[
 2DR\bigl(s_3(\alpha)+s_3(\beta)\bigr)
 =2\bigl(s_3(\alpha)+s_3(\beta)\bigr).
\] Taking absolute values and the fourth-largest entry are both
1-Lipschitz for the maximum norm on these finite lists. Therefore the
truncated section expectation \(\widetilde Q_j\) satisfies \[
 \widetilde Q_j\le Q_j+4\mathbb Es_3
 \le q+4q/D^2=:\widetilde q.                                \tag{5}
\] No regularity or atom assumption is used here.

For each unordered pair \(ij\), let \(E_{ij}\) be the Borel event that
\(i,j\) are the selected two largest coordinates. Define its
\emph{weighted} mass by \[
 m_{ij}=\mathbb E\left[
   \min(|\alpha_i|,|\alpha_j|)\,\mathbf1_{E_{ij}}\right].     \tag{6}
\] Zero second coordinates contribute zero, so arbitrary tie choices
there are harmless. Consider distinct \(i,j,k\). On
\(E_{ij}\times E_{jk}\), the two truncated vectors are exactly supported
on \(ij\) and \(jk\). Put \(r_\alpha=\min(|\alpha_i|,|\alpha_j|)\) and
\(r_\beta=\min(|\beta_j|,|\beta_k|)\). Four entries in (4) have absolute
values \[
 D|\alpha_j|,\quad D|\beta_j|,\quad
 D|\alpha_i\beta_j|,\quad D|\alpha_j\beta_k|.
\] Since \(r_\alpha,r_\beta\le R\) and \(D\le1\), each of these is at
least \(D^2r_\alpha r_\beta\). Integrating the fourth cofactor over this
product event gives \[
 D^2m_{ij}m_{jk}\le\widetilde Q_j\le\widetilde q.           \tag{7}
\] In particular, edges of weight strictly greater than
\(w=\sqrt{\widetilde q}/D\) cannot share a vertex. They form a matching.
Use those edges as pair blocks and every remaining coordinate as a
singleton block. Every discarded edge has weight at most \(w\).

If the selected top pair of \(\alpha\) is a matching edge, its truncated
vector lies in its plane block. Otherwise, its distance to the block
containing its largest coordinate is at most the smaller retained
coordinate. By (3) and the triangle inequality, \[
 \operatorname{dist}(\alpha,\mathcal W)
 \le\sqrt d\,s_3(\alpha)
  +\min(|\alpha_i|,|\alpha_j|)
       \mathbf1_{\{ij\text{ discarded}\}}.
\] Hence (2), (5), and (6) give \[
 \mathbb E\operatorname{dist}(\alpha,\mathcal W)
 \le\frac{\sqrt d\,q}{D^2}
   +\binom d2\frac{\sqrt{q+4q/D^2}}D
 \le\frac{\sqrt d+\sqrt5\binom d2}{D^2}\sqrt q.             \tag{8}
\] The last step uses \(q\le1\) and \(D\le1\). Finally, \(D\ge a/2\),
\(\sqrt5\le5/2\), and \(\sqrt d\le d\) imply \[
 4\left(\sqrt d+\sqrt5\binom d2\right)
 \le4d+5d(d-1)\le5d^2,
\] which proves (1). At \(q=0\), (7) says all positive-weight edges form
a matching, (2) gives exact two-sparsity almost everywhere, and (8) has
zero right-hand side. This proves the zero case directly.

\subsubsection{A sharp obstruction under the intermediate
hypotheses}\label{a-sharp-obstruction-under-the-intermediate-hypotheses}

\textbf{Lemma 2 (cofactor vanishing for a block union).} Let
\(\mathbb R^d=\bigoplus_jW_j\), where \(\dim W_j\in\{1,2\}\). For every
\(d+1\) vectors in \(\bigcup_jW_j\), the fourth absolute cofactor is
zero.

\textbf{Proof.} If the tuple has rank below \(d\), all cofactors vanish.
Otherwise assign each vector to one block containing it, assigning a
zero vector arbitrarily. If block \(j\) receives \(n_j\) vectors, full
rank requires \(n_j\ge\dim W_j\) for every block. As their sum is
\(d+1\), exactly one block receives one extra vector and the others
receive precisely their dimensions. The unique relation on the tuple is
supported on that oversampled block, which receives at most three
vectors. Cofactors are its relation coefficients, so at most three
cofactors are nonzero. \(\square\)

\textbf{Lemma 3 (a quantitative distance witness).} If
\(z_1,\ldots,z_{d+1}\) are in the unit ball, then, for every block union
\(\mathcal V\) in Lemma 2, \[
 \sum_{i=1}^{d+1}\operatorname{dist}(z_i,\mathcal V)
 \ge\phi(z_1,\ldots,z_{d+1}).                               \tag{9}
\]

\textbf{Proof.} Project each \(z_i\) orthogonally onto a closest block
and call the resulting vector \(y_i\). Then \(|y_i|\le1\) and
\(|z_i-y_i|=\operatorname{dist}(z_i,\mathcal V)\). For each cofactor,
multilinear telescoping and Hadamard's inequality bound its change by
the sum of these distances over its \(d\) columns, hence by the full
sum. The fourth absolute cofactor is 1-Lipschitz in the list of
cofactors, and it vanishes for the \(y_i\) by Lemma 2. \(\square\)

Now put \[
 v_{12}=(e_1+e_2)/\sqrt2,\qquad v_{23}=(e_2+e_3)/\sqrt2,
 \qquad 0<p\le1/(d+2),\qquad w=(1-2p)/d.
\] Let \(\nu_p\) give total mass \(w\) to each antipodal pair
\(\{\pm e_i\}\), and mass \(p\) to each of
\(\{\pm v_{12}\},\{\pm v_{23}\}\), distributing each antipodal mass
equally. This is an even law in the unit ball, and \(w\ge p\). Sampling
the \(d\) axis directions in any order gives \[
 A(\nu_p)\ge d!w^d\ge a_d:=d!/(d+2)^d>0.                 \tag{10}
\] Any repeated direction gives a relation supported on its two repeated
columns. If the tuple has rank below \(d\), all cofactors vanish;
otherwise its unique relation is that repeated-direction relation, so at
most two cofactors are nonzero. Thus such a tuple has fourth cofactor
zero. Among distinct-direction \(d+1\)-tuples, the only set with
positive fourth cofactor is \[
 (e_1,e_3,e_4,\ldots,e_d,v_{12},v_{23}).                     \tag{11}
\] Indeed, omitting an outside axis gives rank below \(d\); omitting one
of the two new directions leaves a three-term circuit; and omitting
\(e_1\) or \(e_3\) likewise leaves the three-term circuit on
\(e_2,e_3,v_{23}\) or \(e_1,e_2,v_{12}\). In (11), the four nonzero
absolute cofactors are \(1/2,1/2,1/\sqrt2,1/\sqrt2\), and all other
cofactors are zero. Its fourth cofactor is \(1/2\). Therefore, exactly,
\[
 \varepsilon_p:=\mathbb E_{\nu_p^{d+1}}\phi
 =\frac{(d+1)!}{2}\,w^{d-1}p^2.                            \tag{12}
\] Signs and permutations do not change this value.

Applying Lemma 3 to (11), and using that every direction has mass at
least \(p\), gives for \emph{every} direct sum of line and plane blocks
\[
 \mathbb E_{\nu_p}\operatorname{dist}(X,\mathcal V)
 \ge p/2.                                                  \tag{13}
\] This lower bound is independent of the choice, angles, or uniqueness
of the decomposition. For any basis \(U\) with unit-length columns, let
\(\mathcal V=U\mathcal W\). Since \(\|U\|\le\sqrt d\), \[
 \mathbb E\operatorname{dist}(U^{-1}X,\mathcal W)
 \ge\frac1{\sqrt d}\mathbb E\operatorname{dist}(X,\mathcal V)
 \ge\frac p{2\sqrt d}.                                    \tag{14}
\] As \(p\downarrow0\), (12) is a positive constant depending on \(d\)
times \(p^2+O_d(p^3)\), while (10) is uniformly positive. Thus an
estimate with power \(\gamma>1/2\) and constants depending only on
\(a,d\) is impossible under the hypotheses of Theorem 1.

The obstruction is about arbitrary bounded even laws. No claim is made
that these particular laws are cone laws of centrally symmetric bodies.
The completed cone-specific half-mass argument must still be retained
when applying the matching theorem to convex bodies.

\subsubsection{Consequence for the existing proof
chain}\label{consequence-for-the-existing-proof-chain}

Keeping the completed estimate \(\mathbb E\phi\le F_d\delta^{2/3}\), its
determinant lower bound \(a=a_d^{\mathrm{John}}\), and its geometry
lemma unchanged, (1) gives \[
 s\le5d^2(a_d^{\mathrm{John}})^{-2}
       \sqrt{4(d+1)F_d/a_d^{\mathrm{John}}}\;\delta^{1/3}.
\] The same cap-to-containment step therefore yields a proved upper-end
power \(1/(3d)\), with explicit constants obtained by this substitution.
The root assembly is responsible for stating and simplifying the final
global constant. The intermediate square-root matching step is sharp, so
improving it further requires additional cone-law information or a
different geometric route.

\subsubsection{Verification scope}\label{verification-scope}

\texttt{check\_matching\_improvement.py} checks the cofactor formula,
the truncation perturbation inequality, the adjacent-edge lower bound,
and the exact finite-law obstruction formula using rational arithmetic.
For the last check it uses the rational variant \(v_{ij}=(e_i+e_j)/2\),
whose fourth cofactor is \(1/4\) and whose expectation is
\((d+1)!w^{d-1}p^2/4\), avoiding any floating-point square roots. These
finite checks support the polynomial identities; the measure-theoretic
selection and integral proof above do not depend on enumeration.

\newpage

\section{IV Direct geometric
recovery}\label{iv-direct-geometric-recovery}

\subsection{Direct quantitative recovery from off-block
normals}\label{direct-quantitative-recovery-from-off-block-normals}

Date: 7 October 2026. This is a separate research note. It does not
alter the completed upper-end stability supplement.

\subsubsection{1. Statement}\label{statement-1}

Let \(d\ge2\), and let \(Q\subset\mathbb R^d\) be an origin-symmetric
(\(Q=-Q\)), full-dimensional convex body such that

\[
rB_2^d\subset Q\subset RB_2^d,\qquad 0<r\le R.
\]

Let \(\mathbb R^d=\bigoplus_{j=1}^m W_j\) be an orthogonal decomposition
into nonzero subspaces, let \(\mathcal W=\bigcup_jW_j\), and put

\[
P=\prod_{j=1}^m\operatorname{proj}_{W_j}Q,
\qquad
I=\int_{\partial Q}\operatorname{dist}(n_Q(x),\mathcal W)
       \,d\mathcal H^{d-1}(x).
\]

Regular boundary points suffice in this integral; the omitted set has
surface measure zero. Thus

\[
I=\int_{S^{d-1}}\operatorname{dist}(u,\mathcal W)\,dS_Q(u).
\]

Write \(\kappa_k=|B_2^k|\), including \(\kappa_0=1\), and set

\[
 A_{d,r,R}=\frac{2}{\kappa_{d-2}}
       \left(\frac{2(\sqrt2R+r)}r\right)^{d-2}.
\]

\textbf{Theorem.} The product built from the actual block projections
satisfies

\[
 d_H(Q,P)\le m(A_{d,r,R}I)^{1/(d-1)},                 \tag{1.1}
\]

\[
 d_{\rm BM}(Q,P)-1
 \le\frac mr(A_{d,r,R}I)^{1/(d-1)}.                 \tag{1.2}
\]

Equivalently, for the cone-law measurement

\[
s=\mathbb E_{\nu_Q}\operatorname{dist}(X,\mathcal W)
  =\frac I{d|Q|},
\]

one has the explicit estimate

\[
\boxed{d_{\rm BM}(Q,P)-1
\le \frac mr
\left(A_{d,r,R}\,d\kappa_d R^d\,s\right)^{1/(d-1)}.}  \tag{1.3}
\]

No atomicity, polyhedrality, smoothness, or uniqueness of a product
decomposition is required. The dimensions of the blocks do not enter the
proof. If every block has dimension one or two, \(P\) is in the equality
class of the completed theorem, so its distance to the entire class has
the same upper bound. For \(m=1\), \(P=Q\), and all claims are
immediate.

\subsubsection{2. A planar mixed-normal
lemma}\label{a-planar-mixed-normal-lemma}

\textbf{Lemma.} Let \(C\subset\mathbb R^2\) be a full-dimensional
compact convex set. Suppose \(C\) contains points \((-b_0,0)\) and
\((0,-a_0)\), where \(a_0,b_0>0\), while
\(\operatorname{dist}(0,C)\ge t>0\). Then

\[
\int_{\partial C}\min(|n_1|,|n_2|)\,d\mathcal H^1\ge t. \tag{2.1}
\]

\textbf{Proof.} Intersect with the closed negative quadrant. Choose the
closest intersection points to the origin on its coordinate axes and
write them as \((-b,0)\) and \((0,-a)\), with \(a,b>0\). They exist
because the original two points are in \(C\), and positivity follows
from the distance assumption. For \(-b\le x\le0\), let

\[
 f(x)=\max\{y:(x,y)\in C\cap(-\infty,0]^2\}.
\]

The segment connecting the two chosen points proves that this function
is defined on the entire interval. Convexity shows that \(f\) is
concave; it is continuous, \(f(-b)=0\), and \(f(0)=-a\). The choice of
\(b\) gives \(f(x)<0\) for \(-b<x\le0\). Concavity and \(f\le0\) imply
that \(f\) is nonincreasing. Its open graph belongs to the boundary of
the original \(C\), because both quadrant constraints are inactive
there.

Put \(g=-f\). This is a nonnegative, nondecreasing convex function, with
\(g(-b)=0\) and \(g(0)=a\). It is absolutely continuous and its
almost-everywhere derivative is nonnegative and nondecreasing. At
regular graph points the unit outward normal is
\((g'(x),1)/\sqrt{1+g'(x)^2}\), so the graph contribution is

\[
\int_{-b}^0\min(1,g'(x))\,dx.
\]

Since \(g'\) is nondecreasing, let \(x_*\in[-b,0]\) be a crossing point
of the level one, with the natural endpoint choices if \(g'\ge1\)
everywhere or \(g'\le1\) everywhere. Integration before and after the
crossing gives

\[
\int_{-b}^0\min(1,g')\,dx=g(x_*)-x_*
                      =\min_{[-b,0]}(g(x)-x).
\]

Each graph point \((x,-g(x))\in C\) is at Euclidean distance at least

\(t\) from the origin. Its distance in the \(\ell_1\) norm is
\(g(x)-x\), hence the displayed minimum is at least \(t\). All other
boundary contributions are nonnegative. This proves (2.1). Endpoint
vertical segments, if present, have zero mixed-normal contribution and
do not change the argument. \(\square\)

\subsubsection{3. Two-factor recovery by thickening a planar
section}\label{two-factor-recovery-by-thickening-a-planar-section}

Fix an orthogonal bipartition \(\mathbb R^d=E\oplus F\) into nonzero
subspaces and put

\[
 P_{E,F}=\operatorname{proj}_E Q\times\operatorname{proj}_F Q,
\qquad
 I_{E,F}=\int_{\partial Q}
        \min(|\operatorname{proj}_E n|,|\operatorname{proj}_F n|)
        \,d\mathcal H^{d-1}.
\]

We claim

\[
 d_H(Q,P_{E,F})^{d-1}\le A_{d,r,R}I_{E,F}.            \tag{3.1}
\]

Since \(Q\subset P_{E,F}\), choose a point \(p\in P_{E,F}\) where
\(t=\operatorname{dist}(p,Q)=d_H(Q,P_{E,F})\) is attained. If \(t=0\)
the claim is immediate. Choose \(q_A,q_B\in Q\) with

\[
 \operatorname{proj}_E q_A=\operatorname{proj}_E p,
\qquad
 \operatorname{proj}_F q_B=\operatorname{proj}_F p.
\]

Then \(a=p-q_A\in F\) and \(b=p-q_B\in E\) are nonzero and orthogonal.
Define the two-dimensional plane

\[
 L=\operatorname{span}\{b,a\},\quad
 e_1=b/|b|\in E,\quad e_2=a/|a|\in F,\quad T=L^\perp.
\]

Because each component of \(p\) is a projection of a point of \(Q\),
\(|p|\le\sqrt2R\). Also \(t\le |p|\le\sqrt2R\), because \(0\in Q\). Set

\[
 \eta=\frac{t}{2(\sqrt2R+r)}\in(0,1).
\]

For every \(z\in T\) with \(|z|<\eta r\), the points

\[
 q_{A,z}=(1-\eta)q_A+z,\qquad
 q_{B,z}=(1-\eta)q_B+z
\]

belong to the interior of \(Q\): each is the convex combination of a
point of \(Q\) and the interior-ball point \(z/\eta\). They share an
affine plane parallel to \(L\) with

\[
 p_z=(1-\eta)p+z.
\]

Relative to \(p_z\), the endpoints have coordinates \((0,-(1-\eta)|a|)\)
and \((-(1-\eta)|b|,0)\). Furthermore

\[
 \operatorname{dist}(p_z,Q)
 \ge t-|p-p_z|
 \ge t-\eta\sqrt2R-|z|
 \ge t/2.                                           \tag{3.2}
\]

The planar section \(C_z=Q\cap(p_z+L)\) is full-dimensional in its plane
and contains both endpoints. Applying the planar lemma gives

\[
 \int_{\partial C_z}
       \min(|n_{C_z}\cdot e_1|,|n_{C_z}\cdot e_2|)
       \,d\mathcal H^1\ge t/2.                     \tag{3.3}
\]

The linear projection \(\pi_T:\partial Q\to T\) has tangential
\((d-2)\)-Jacobian \(|\operatorname{proj}_L n|\). For almost every
transverse parameter, its slice has outward normal

\[
 n_{C_z}=\frac{\operatorname{proj}_L n}
                 {|\operatorname{proj}_L n|}
\]

at almost every regular intersection point. The coarea formula therefore
gives, with zero integrand when \(\operatorname{proj}_L n=0\),

\[
\begin{aligned}
I_{E,F}
&\ge\int_{\partial Q}\min(|n\cdot e_1|,|n\cdot e_2|)
                       \,d\mathcal H^{d-1}\\
&\ge\int_{\{|z|<\eta r\}\subset T}
   \int_{\partial C_z}
       \min(|n_{C_z}\cdot e_1|,|n_{C_z}\cdot e_2|)
       \,d\mathcal H^1\,dz\\
&\ge\kappa_{d-2}(\eta r)^{d-2}\,t/2\\
&=\frac{\kappa_{d-2}}2
       \left(\frac{r}{2(\sqrt2R+r)}\right)^{d-2}t^{d-1}.
\end{aligned}
\]

This is (3.1). For \(d=2\), the transverse space has dimension zero, and
the same formula is interpreted with \(\kappa_0=1\).

The coarea argument applies to arbitrary convex bodies: their boundaries
are countably rectifiable and regular almost everywhere. The singular
boundary has zero \((d-1)\)-measure and therefore zero weighted coarea
contribution. For almost every transverse parameter, the intersection of
the regular boundary is the section boundary up to length-zero sets. The
ball thickening ensures that all sections used here have nonempty
relative interior. Flat facets and nonatomic surface-area laws are both
allowed.

\subsubsection{4. Recovering all factors}\label{recovering-all-factors}

For each \(j\), apply (3.1) to \(E=W_j\), \(F=W_j^\perp\), and write the
resulting two-factor product as \(P_j\). For any unit vector \(n\),

\[
 \min(|\operatorname{proj}_{W_j}n|,
       |\operatorname{proj}_{W_j^\perp}n|)
 \le \operatorname{dist}(n,\mathcal W).
\]

Indeed, if the largest block component is the \(j\)-th one, the
right-hand side is exactly the complementary norm; if another block
component is largest, the right-hand side is at least the \(j\)-th
component norm. Consequently

\[
 e_j:=d_H(Q,P_j)\le(A_{d,r,R}I)^{1/(d-1)}.            \tag{4.1}
\]

Fix any \(p\in P\), start with \(q_0\in Q\), and successively replace
the \(j\)-th block coordinate of \(q_{j-1}\) by that of \(p\). Call the
new point \(z_j\). It belongs to \(P_j\), because its \(j\)-th component
is feasible in the projection of \(Q\), while the complementary
component is the projection of \(q_{j-1}\in Q\). Choose \(q_j\in Q\)
such that

\[
 q_j=z_j+\varepsilon_j,\qquad |\varepsilon_j|\le e_j.
\]

Writing \(\pi_j=\operatorname{proj}_{W_j}\), the recurrence is

\[
 q_j=\pi_jp+(I-\pi_j)q_{j-1}+\varepsilon_j.
\]

The mutually orthogonal coordinate projections commute. After all blocks
have been replaced, the original \(q_0\) contributes zero and

\[
 q_m-p=\sum_{j=1}^m
       \left(\prod_{k=j+1}^m(I-\pi_k)\right)\varepsilon_j.
\]

Every displayed product is an orthogonal projection, so
\(|q_m-p|\le\sum_j e_j\). Taking the maximum over \(p\in P\) proves
(1.1).

Finally, if \(e=d_H(Q,P)\), support functions satisfy

\[
 h_P(u)\le h_Q(u)+e\le(1+e/r)h_Q(u)
\]

for unit \(u\), because \(h_Q(u)\ge r\). Hence

\[
 (1+e/r)^{-1}P\subset Q\subset P.
\]

This proves (1.2), and \(|Q|\le\kappa_dR^d\) gives (1.3).

\subsubsection{5. Consequences and
limits}\label{consequences-and-limits}

The weighted corner cuts in \texttt{geometry-next.md}, and their
all-product distance obstruction imported from the audited supplement,
show that the exponent \(1/(d-1)\) in (1.3) cannot be improved in
general when the blocks have dimension one or two. Their radii can be
kept uniformly bounded above and below. Thus this direct recovery
theorem closes the geometry-only exponent gap; it does not identify the
sharp exponent of the original upper-end deficit.

The selected basis in the completed rank lemma already gives an explicit
inradius sufficient to apply this theorem. If its columns have Euclidean
norm at most one and determinant \(D\ge a_d/2\), every row of its
inverse has norm at most \(1/D\) by the cofactor-wedge bound. Therefore

\[
 \|U^{-1}\|\le\sqrt d/D,\qquad
 U^TK\supset \frac{a_d}{2\sqrt d}B_2^d
\]

when \(K\) is in symmetric John position. The same transformed body
satisfies \(U^TK\subset dB_2^d\). Thus the existing estimate
\(s\le L_d\delta^{1/6}\) yields an improved geometric exponent
\(\delta^{1/[6(d-1)]}\), with constants obtained directly from (1.3).
The direct geometric estimate and its use in the stronger matching
assembly passed an independent analytic model audit. It does not change
the completed release and does not rely on the extra-factor estimate
refuted in \texttt{geometry-next.md}.

\subsubsection{6. Stronger homothetic recovery without an inradius
constant}\label{stronger-homothetic-recovery-without-an-inradius-constant}

The interior-ball thickening can be replaced by a thickening with \(Q\)
itself. This removes the inradius from the small-error bound.

Keep \(Q=-Q\), \(Q\subset RB_2^d\), and the orthogonal block
decomposition from above. Define

\[
T=\left(2^{2d-1}\,dR\,s\right)^{1/(d-1)},\qquad
\lambda=\max\{\alpha\ge0:\alpha P\subset Q\}.
\]

\textbf{Theorem.} For arbitrary nonzero block dimensions,

\[
\boxed{1-\lambda\le mT.}                              \tag{6.1}
\]

If \(mT\le1/2\), then

\[
d_{\rm BM}(Q,P)-1\le2mT.                              \tag{6.2}
\]

When all block dimensions are one or two, the global whole-class
estimate is

\[
\boxed{D(Q,\mathcal E_d)-1
       \le2md\left(2^{2d-1}dR\,s\right)^{1/(d-1)}.}    \tag{6.3}
\]

The global estimate uses the symmetric John bound only for the regime
\(mT>1/2\). It is an estimate to the entire affine product class; in
that regime it does not claim control of distance to the one fixed
product \(P\).

\paragraph{6.1. Bipartition lemma}\label{bipartition-lemma}

For an orthogonal bipartition \(E\oplus F\), put \(P_{E,F}\) and
\(I_{E,F}\) as in Section 3, and set

\[
\lambda_{E,F}=\max\{\alpha:\alpha P_{E,F}\subset Q\},
\qquad \tau=1-\lambda_{E,F}.
\]

We prove

\[
\boxed{\tau^{d-1}\le
       2^{2d-1}R\,\frac{I_{E,F}}{|Q|}.}                \tag{6.4}
\]

Both bodies contain the origin in their interiors, and \(Q\subset
P_{E,F}\), so \(\lambda_{E,F}\in(0,1]\). The claim is immediate when
\(\tau=0\). Otherwise compactness of the support-function ratio gives a
unit vector \(u\) such that

\[
h_Q(u)=\lambda_{E,F}h_{P_{E,F}}(u).
\]

Write \(H=h_{P_{E,F}}(u)>0\), and choose \(p\in P_{E,F}\) with
\(u\cdot p=H\). Choose \(q_A,q_B\in Q\) sharing respectively the \(E\)
and \(F\) components of \(p\). As before,

\[
a=p-q_A\in F,\quad b=p-q_B\in E,\quad
L=\operatorname{span}\{b,a\},\quad T_0=L^\perp.
\]

The two vectors \(a,b\) are nonzero: the support gap at \(p\) is
\(\tau H>0\). Furthermore \(u_L=\operatorname{proj}_L u\ne0\), because

\[
u\cdot a=H-u\cdot q_A\ge\tau H>0.
\]

Set \(\eta=\tau/4\in(0,1/4)\). For any
\(z\in\operatorname{int}_{T_0}(\operatorname{proj}_{T_0}Q)\), choose an
interior lift \(q\in\operatorname{int}Q\) with
\(\operatorname{proj}_{T_0}q=z\). Such a lift exists because a
surjective linear map sends the interior of a convex body onto the
interior of its image. Consider

\[
p_q=(1-\eta)p+\eta q,\quad
q_{A,q}=(1-\eta)q_A+\eta q,\quad
q_{B,q}=(1-\eta)q_B+\eta q.
\]

The endpoints belong to the interior of \(Q\), share a plane parallel to
\(L\) with \(p_q\), and relative to \(p_q\) have the same negative-axis
form as in Section 3. Central symmetry gives \(u\cdot q\ge-h_Q(u)\).
Consequently

\[
\begin{aligned}
u\cdot p_q-h_Q(u)
 &=\tau H-\eta(H-u\cdot q)\\
 &\ge\tau H-\eta(H+h_Q(u))\\
 &\ge\tau H-2\eta H=\tau H/2.
\end{aligned}
\]

On this fixed affine plane, only the component \(u_L\) of the normal
varies. Thus the distance of \(p_q\) to the planar section is at least

\[
\frac{\tau H}{2|u_L|}.
\]

Apply the planar lemma and then the same weighted coarea identity as in
Section 3. The transverse parameters now fill

\[
(1-\eta)\operatorname{proj}_{T_0}p+
       \eta\operatorname{proj}_{T_0}Q,
\]

whose \((d-2)\)-volume is \(\eta^{d-2}|\operatorname{proj}_{T_0}Q|\).
The boundary of this projection has measure zero. Every interior
parameter has an interior lift, so its section is full-dimensional; no
measurable selection of the lift is needed, since the lower bound holds
for each such section. It follows that

\[
I_{E,F}\ge
\eta^{d-2}|\operatorname{proj}_{T_0}Q|\,
       \frac{\tau H}{2|u_L|}.                        \tag{6.5}
\]

For any fixed transverse parameter \(w\), its fiber in \(Q\) is
contained in a planar rectangle with side lengths at most

\[
\frac{2h_Q(u)}{|u_L|}\le\frac{2H}{|u_L|}
\quad\hbox{and}\quad 2R.
\]

The first follows from
\(-h_Q(u)\le u_L\cdot x_L+u_{T_0}\cdot w\le h_Q(u)\); the second follows
from \(Q\subset RB_2^d\) in the perpendicular direction inside \(L\).
Fubini therefore yields

\[
|Q|\le\frac{4RH}{|u_L|}\,
          |\operatorname{proj}_{T_0}Q|.              \tag{6.6}
\]

Substitute (6.6) into (6.5), and use \(\eta=\tau/4\):

\[
\frac{I_{E,F}}{|Q|}
\ge\frac{\tau^{d-1}}{8\,4^{d-2}R}
 =\frac{\tau^{d-1}}{2^{2d-1}R}.
\]

This proves (6.4), also in dimension two with the zero-dimensional
transverse volume convention.

\paragraph{6.2. Combining the homothetic
inclusions}\label{combining-the-homothetic-inclusions}

For each single-block bipartition \(P_j=\operatorname{proj}_{W_j}Q\times
\operatorname{proj}_{W_j^\perp}Q\), let
\(\lambda_j=\max\{\alpha:\alpha P_j\subset Q\}\). The pointwise normal
inequality from Section 4 and \(I=d|Q|s\) give

\[
1-\lambda_j\le T.                                   \tag{6.7}
\]

We claim

\[
\left(\prod_{j=1}^m\lambda_j\right)P\subset Q.         \tag{6.8}
\]

Fix \(p\in P\), write \(p_j=\operatorname{proj}_{W_j}p\), and prescribe
the adjusted coordinates

\[
\widetilde p_j=\left(\prod_{k<j}\lambda_k\right)p_j.
\]

Each adjusted coordinate lies in its factor projection, since that
factor is convex and contains zero. Start with \(q_0=0\in Q\). For
\(j=1,\ldots,m\), replace the \(j\)-th block of \(q_{j-1}\) by
\(\widetilde p_j\) to obtain \(z_j\in P_j\), and then set
\(q_j=\lambda_j z_j\in Q\). The \(j\)-th block of \(q_m\) is

\[
\left(\prod_{k=j}^m\lambda_k\right)\widetilde p_j
=\left(\prod_{k=1}^m\lambda_k\right)p_j.
\]

Therefore \(q_m=(\prod_j\lambda_j)p\in Q\), proving (6.8). Since the
\(\lambda_j\) lie in \((0,1]\),

\[
1-\prod_j\lambda_j\le\sum_j(1-\lambda_j)\le mT.
\]

This proves (6.1). In the small-error regime it gives

\[
(1-mT)P\subset Q\subset P,\qquad
d_{\rm BM}(Q,P)-1
 \le\frac{mT}{1-mT}\le2mT,
\]

which is (6.2). For the whole-class estimate, the same bound applies
when \(mT\le1/2\), and otherwise symmetric John's theorem gives
\(D(Q,\mathcal E_d)-1\le d-1\le2mdT\). This proves (6.3).

\paragraph{6.3. Original endpoint
implication}\label{original-endpoint-implication}

The completed proof's final coordinate body satisfies
\(Q\subset dB_2^d\). Thus (6.2) or (6.3) can be used without paying for
the small determinant of the selected basis. In particular, the existing
\(s\le L_d\delta^{1/6}\) improves the deficit exponent from \(1/(6d)\)
to \(1/[6(d-1)]\). Any independent improvement in the matching-stage
power can be combined with this proof, after checking its precise
inputs.

The \(1/(d-1)\) power in (6.3) is sharp in fixed dimension by the corner
cut bodies and the all-product distance lower bound in geometry-next.md.
The new proof does not infer a sharp original-deficit exponent from that
intermediate sharpness.

The direct recovery proof uses thickening by the body itself and a
shared support direction to cancel the conditioning loss. Its
small-error estimate needs no inradius constant and has geometric power
\(1/(d-1)\). The optimal original-deficit exponent remains unresolved.

\newpage

\section{V Half mass obstruction for actual cone
laws}\label{v-half-mass-obstruction-for-actual-cone-laws}

\subsection{Sharpness of the cone-specific half-mass
gate}\label{sharpness-of-the-cone-specific-half-mass-gate}

7 October 2026. This note concerns the short-cofactor gate in Sections
2--3 of the completed upper-end stability manuscript. It does not modify
that manuscript or assert an improved stability exponent.

\subsubsection{Result}\label{result}

For every fixed dimension \(d\ge3\), the estimate \[
 \mathbb E\phi\le F_d\delta^{2/3}
\] has the best possible power under its stated hypotheses, even among
centrally symmetric polytopes already in symmetric John position. Here
\(\delta=1/2-a(K)\), and \(\phi\) is the fourth-largest absolute
horizontal cofactor of \(d+1\) independent cone-law samples. Neither the
corner-slab area power two nor the weighted near-half-mass tail power
two can be improved without adding hypotheses or changing the
intermediate quantity.

The obstruction is an actual cone law. It does not use a freely
prescribed even law on a polar boundary.

\subsubsection{1. A product of a truncated three-cube with
intervals}\label{a-product-of-a-truncated-three-cube-with-intervals}

For \(0<t\le1\), define \[
 H_t=[-1,1]^3\cap\{|x_1+x_2+x_3|\le3-t\},\qquad
 K_{d,t}=H_t\times[-1,1]^{d-3}.
\] The interpretation in dimension three is \(K_{3,t}=H_t\). The unit
ball is contained in \(K_{d,t}\): its first-three-coordinate sum has
absolute value at most \(\sqrt3<2\le3-t\). The body is contained in the
cube and hence in \(\sqrt d B_2^d\). Its coordinate-facet contacts
\(\pm e_i\), with weights \(1/2\) for each sign, give the symmetric John
identity \(\sum_i(e_ie_i^{\mathsf T})=I\). Thus the unit ball is its
John ellipsoid; there is no hidden change of normalization in the
cofactor calculation.

Put \[
 Z=24-t^3,\qquad p=\frac{8-t^2}{Z},\qquad
 r=\frac{(3-t)t^2}{Z},\qquad q=\frac1{3-t}.
\] The cone law of \(H_t\) has three direction pairs \(\pm e_i\), each
with total probability \(p\), and one pair \(\pm q(e_1+e_2+e_3)\), with
total probability \(r\). Signs within each pair are uniform. In
particular \(3p+r=1\).

Indeed, \(|H_t|=8-t^3/3\). Each coordinate facet has area \(4-t^2/2\)
and support height one, giving pair cone mass \(8-t^2\). Each new
triangular facet has area \(\sqrt3t^2/2\), support height
\((3-t)/\sqrt3\), and scaled normal \(\pm q(e_1+e_2+e_3)\). Its two-sign
pair cone mass is \((3-t)t^2\). The normalization is \(3|H_t|=Z\).

The product cone law is supported on the same four direction pairs in
the first three coordinates, together with \(\pm e_4,\ldots,\pm e_d\).
Its pair probabilities are \[\begin{gathered}
P=\frac3d p\quad\text{for each of the first three axes},\\
 R=\frac3d r\quad\text{for the diagonal pair},\\
 w=\frac1d\quad\text{for every remaining axis}.
\end{gathered}\] To verify the dimension factors directly, a facet lifted from the
three-dimensional factor has its area-times-height and total body volume
both multiplied by \(2^{d-3}\); division by \(d|K_{d,t}|\) instead of
\(3|H_t|\) multiplies the corresponding normalized mass by \(3/d\).
Every interval-axis pair has area-times-height \(|K_{d,t}|\), and
therefore normalized mass \(1/d\). This argument uses ordinary facet
cone-volume mass and does not assume that an arbitrary discrete measure
is realizable.

\subsubsection{2. Exact fourth-cofactor
expectation}\label{exact-fourth-cofactor-expectation}

Call the unsigned direction pairs types. The \(d+1\) types span
dimension \(d\), and their unique linear dependence is supported on the
first three axes and the diagonal. A \(d+1\)-tuple with fewer than \(d\)
distinct types has all horizontal cofactors zero. A tuple with exactly
\(d\) distinct types has one repeated parallel direction. If its rank is
\(d\), its kernel is generated by the relation between those two columns
and has just two nonzero cofactors; if its rank is smaller, all
cofactors vanish.

Consequently \(\phi>0\) occurs exactly when all \(d+1\) types occur. On
that event the absolute cofactors are \[
 1,\ q,\ q,\ q,\ \underbrace{0,\ldots,0}_{d-3},
\] regardless of the independent signs. Since \(q\le1/2\), their
fourth-largest member is \(q\). The event probability is \[
 \Pi=(d+1)!P^3Rw^{d-3}
      =\frac{81(d+1)!}{d^{d+1}}p^3r.
\] Therefore \[
 \boxed{\mathbb E\phi
  =\frac{81(d+1)!}{d^{d+1}}
    \frac{(8-t^2)^3t^2}{(24-t^3)^4}
  \sim\frac{(d+1)!}{8d^{d+1}}t^2.}                  \tag{2.1}
\]

\subsubsection{3. Exact deficit}\label{exact-deficit}

For completeness this deficit can be computed without invoking the
dimension-weighted product identity. The determinant expectation is \[
 A=d!w^{d-3}P^2(P+3Rq)
   =\frac{27d!}{d^d}p^2(p+3rq).
\] The first term uses all \(d\) axis types; the three remaining nonzero
terms replace one of the first three axes by the diagonal. Replacing any
other axis by the diagonal gives zero determinant.

Tuples with fewer than \(d\) types have zero cofactor mass. Tuples with
exactly \(d\) types have zero balanced Rademacher defect, because their
two nonzero cofactor magnitudes agree. On the all-distinct event, sign
averaging gives \[
 \mathbb E_\epsilon|\epsilon_0+q(\epsilon_1+\epsilon_2+\epsilon_3)|
   =1+\frac14(3q-1).
\] Here only the two all-equal patterns in the last three signs have
\(|q(\epsilon_1+\epsilon_2+\epsilon_3)|>1\); all other patterns have
value \(q<1\). Its unnormalized cofactor defect is consequently \[
 \frac{1+3q}{2}-\left(1+\frac14(3q-1)\right)
   =\frac{3q-1}{4}=\frac{t}{4(3-t)}.
\] Using \(\Delta=(d+1)A\delta\), this proves \[
 \boxed{\delta_{d,t}
    =\frac3d\frac{t^3(8-t^2)}{8(24-t^3)(4+t^2)}
    \sim\frac1{32d}t^3.}                                  \tag{3.1}
\] In particular the normalization agrees with the completed
manuscript's dimension-three formula.

Combining (2.1) and (3.1), for every \(\beta>2/3\), \[
 \frac{\mathbb E\phi}{\delta_{d,t}^{\beta}}
 \sim\frac{(d+1)!}{8d^{d+1}}(32d)^\beta t^{2-3\beta}
 \longrightarrow\infty.
\] No dimension-dependent finite constant can repair an exponent
\(\beta>2/3\) in this intermediate estimate, in any fixed \(d\ge3\).

\subsubsection{4. Sharpness of the weighted tail and contact-slab
powers}\label{sharpness-of-the-weighted-tail-and-contact-slab-powers}

On the all-distinct event, \[
 M=1+3q=\frac{6-t}{3-t},\quad
 t_4=\frac1{6-t},\quad
 g=\frac12-\frac1M=\frac{t}{2(6-t)}.
\] All other tuples have \(t_4=0\). Choose \(\rho=1/7\) and
\(\eta_t=t/[2(6-t)]\). Then \(t_4>\rho\) on every all-distinct tuple and
\(\eta_t\le1/10<1/4\). Thus \[
 \mathbb E\bigl[M\mathbf1_{\{t_4>\rho,\ g\le\eta_t\}}\bigr]
  =M\Pi
  =\frac{81(d+1)!}{d^{d+1}}
    \frac{(8-t^2)^3t^2(6-t)}{(24-t^3)^4}
  \sim\frac{3(d+1)!}{4d^{d+1}}t^2.
\] Since \(t=12\eta_t/(1+2\eta_t)\), \[
 \lim_{t\downarrow0}
 \frac{\mathbb E[M\mathbf1_{\{t_4>1/7,\ g\le\eta_t\}}]}
      {\eta_t^2}
  =\frac{108(d+1)!}{d^{d+1}}>0.                       \tag{4.1}
\] Therefore the exponent two in the weighted near-half-mass bound
cannot be increased, even with \(\rho\) fixed and even for these actual
cone laws.

The corner-slab area estimate is independently sharp. The two diagonal
facets of \(K_{d,t}\) lie in the respective three-coordinate corner
slabs of width \(t\): if \(x_1+x_2+x_3=3-t\), then each \(x_i\ge1-t\),
with the analogous negative assertion. Their total regular boundary area
is \[
 \sqrt3\,2^{d-3}t^2.
\] Any area bound with a power strictly greater than two would fail as
\(t\downarrow0\). Flat facets, rather than exceptional nonregular
contacts, already realize this obstruction. The completed
projection-area and determinant-cancellation argument therefore cannot
acquire a better power merely by a stronger universal patch-area
estimate.

\subsubsection{5. A different geometric quantity does avoid this loss on
the
family}\label{a-different-geometric-quantity-does-avoid-this-loss-on-the-family}

The preceding obstruction concerns an intermediate moment. It does not
exclude a stronger final stability exponent through another argument.
Let \(C_d=[-1,1]^d\) and consider the actual cone support surplus \[
 v=\mathbb E_{\nu_{K_{d,t}}}[h_{C_d}(X)-1].
\] Axis samples have zero surplus. At a diagonal sample the surplus is
\(3q-1=t/(3-t)\), and its probability is \(R\). Consequently \[
 \boxed{v=\frac{3t^3}{d(24-t^3)},\qquad
   \frac v{\delta_{d,t}}=\frac{8(4+t^2)}{8-t^2}
     \le\frac{40}{7}.}                                  \tag{5.1}
\] The actual mixed-volume error used in product recovery is therefore
linear in \(\delta\) on this sharp half-mass family. Bounding the
support surplus by a support-function Lipschitz constant times distance
to the coordinate line/plane union discards a factor of order \(t\)
here: the diagonal samples stay a positive distance from that union,
while their support surplus tends to zero. A proof of a global joint
block-selection and support-surplus estimate would be a different route;
this note proves no such estimate.

One loss-free weighted quantity is already available for arbitrary
balanced coefficients. Since \(t_4\le1/4\) and \(g\le1/2\), \[
 t_4g\le\min(t_4/2,g/4)\le\mathscr D(t).
\] Hence the exact deficit identity gives \[
 \boxed{\mathbb E[\phi g]\le\Delta=(d+1)A\delta.}        \tag{5.2}
\] On the truncation family this weighted moment has order \(t^3\),
rather than the order \(t^2\) of \(\mathbb E\phi\). Equation (5.2) is
not by itself a geometric stability theorem: the missing step is to
select one global line/plane decomposition while retaining an actual
support-surplus weight instead of unweighted sparsity distance.

\subsubsection{Verification scope}\label{verification-scope-1}

The analytic argument above proves the all-dimensional statements. An
optional exact-rational regression in
\texttt{check\_halfmass\_sharpness.py} checks the formulas by direct
determinant/type and sign enumeration in dimensions three and four. It
does not replace the proof and imposes no polytope restriction on the
completed upper stability theorem.

\newpage

\section{VI Weighted cuts and geometric
obstruction}\label{vi-weighted-cuts-and-geometric-obstruction}

\subsection{Geometry gate: exact obstructions and one restricted
improvement}\label{geometry-gate-exact-obstructions-and-one-restricted-improvement}

Date: 7 October 2026. This note leaves the completed quantitative
theorem unchanged. It isolates what can and cannot be inferred from the
inputs of its product-recovery lemma. Every assertion below is proved
explicitly; the general recovery theorem subsequently obtained is in the
separate note \texttt{geometry-rigidity.md}.

\subsubsection{1. Measurements and a cancellation worth
preserving}\label{measurements-and-a-cancellation-worth-preserving}

Let \(Q\) be a centered full-dimensional convex body, let
\(\mathbb R^d=\bigoplus_j W_j\) be an orthogonal decomposition, and
write \(\mathcal W=\bigcup_jW_j\),
\(P=\prod_j\operatorname{proj}_{W_j}Q\). Set

\[
s(Q,\mathcal W)=\mathbb E_{\nu_Q}\operatorname{dist}(X,\mathcal W),
\qquad
v(Q,P)=\mathbb E_{\nu_Q}[h_P(X)-1].
\]

Because distance to the union of linear subspaces is homogeneous, the
cone weight cancels exactly:

\[
s(Q,\mathcal W)=\frac1{d|Q|}\int_{S^{d-1}}
       \operatorname{dist}(u,\mathcal W)\,dS_Q(u).       \tag{G1}
\]

There is no support-height factor left in (G1). This is valid for
arbitrary surface-area measures, including atoms. The mixed-volume
identity likewise is

\[
v(Q,P)=\frac1{d|Q|}\int(h_P(u)-h_Q(u))\,dS_Q(u)
       =\frac{V(Q[d-1],P)}{|Q|}-1.                    \tag{G2}
\]

The completed proof bounds \(v\le R(\sqrt m+1)s\), then converts
relative volume loss into a containment loss. The example below
distinguishes those two steps.

\subsubsection{2. A complete exact calculation for weighted corner
cuts}\label{a-complete-exact-calculation-for-weighted-corner-cuts}

Fix \(d\ge3\), positive coefficients \(a_i\le1\) with \(\max_i a_i=1\),
and \(0<t\le\min_i a_i\). Let

\[
S=\sum_i a_i,\quad H=S-t,\quad A_* =\prod_i a_i,\qquad
Q_{a,t}=[-1,1]^d\cap\{x:|\langle a,x\rangle|\le H\}.      \tag{G3}
\]

Both removed corner simplices are contained in their respective cube
corner regions: writing \(y_i=1-x_i\), the positive cap is
\(y_i\ge0,\ \sum_i a_i y_i<t\), and its intercepts are \(t/a_i\le1\).
Their interiors are disjoint. In particular

\[
V:=|Q_{a,t}|=2^d-\frac{2t^d}{d!A_*}.                      \tag{G4}
\]

For any coordinate block of size at most two, its projection is the
entire coordinate cube. Indeed, at any vertex of that block choose signs
among the remaining coordinates so that the full vertex is neither of
the two all-equal vertices. Such a full vertex has
\(|\langle a,x\rangle|\le S-2\min_i a_i\le S-t\). All coordinate-block
cube vertices therefore lift to points of \(Q_{a,t}\), and convexity
proves the claim. Hence \(P=[-1,1]^d\) for every such coordinate
partition. (The same statement holds for a single coordinate projection
without needing a partition.)

Each coordinate facet has area

\[
\alpha_i=2^{d-1}-\frac{t^{d-1}}{(d-1)!\prod_{k\ne i}a_k}.
\]

The two new facet area-normal vectors are

\[
\pm c a,\qquad c=\frac{t^{d-1}}{(d-1)!A_*}.                \tag{G5}
\]

For example, projecting a new facet onto \(e_i^\perp\) gives a
\((d-1)\)-simplex of volume \(t^{d-1}/((d-1)!\prod_{k\ne i}a_k)=ca_i\),
proving (G5) by the projection area formula. All old facet normals
belong to \(\mathcal W\), and the new cone-law sample is \(\pm a/H\).
Consequently the exact formulas are

\[
\boxed{
s(Q_{a,t},\mathcal W)
 =\frac{2t^{d-1}\operatorname{dist}(a,\mathcal W)}
        {d(d-1)!A_* V},\qquad
v(Q_{a,t},P)=\frac{2t^d}{d(d-1)!A_* V}.
}                                                        \tag{G6}
\]

Their ratio is particularly informative:

\[
\frac{v}{s}=\frac{t}{\operatorname{dist}(a,\mathcal W)}.     \tag{G7}
\]

These formulas account for every facet, rather than just the new faces.
The normalization can also be checked directly:

\[
2\sum_i\alpha_i+2cH=dV.
\]

The maximal centered homothetic containment is exactly

\[
\lambda P\subseteq Q_{a,t},\qquad
\lambda=\frac{S-t}{S},\qquad \tau:=1-\lambda=\frac tS.     \tag{G8}
\]

Indeed, the sole new slab constraint forces and suffices for this
factor.

\subsubsection{3. Two sharp upper restrictions, with different
exponents}\label{two-sharp-upper-restrictions-with-different-exponents}

For \(a=(1,\ldots,1)\), write \(K_t=Q_{a,t}\). If the coordinate
partition has largest block dimension \(b\in\{1,2\}\), then

\[
s(K_t,\mathcal W)=
\frac{2\sqrt{d-b}}{d(d-1)!|K_t|}\,t^{d-1},\qquad
v(K_t,P)=\frac{2}{d(d-1)!|K_t|}\,t^d.                     \tag{G9}
\]

Here \(P=[-1,1]^d\), \(\tau=t/d\), and
\(\theta=1-|K_t|/|P|=2t^d/(d!2^d)\). Thus the
relative-volume-to-homothetic-containment exponent \(1/d\) is optimal
already in this nested symmetric family. It cannot be increased even
when both bodies have uniformly bounded radii and a fixed positive
inradius.

For the stronger conclusion involving distance to the \textbf{entire}
product class, the completed and independently audited Theorem B proves

\[
D(K_t,\mathcal E_d)-1\ge\frac{t}{5d^2}.                    \tag{G10}
\]

Combining (G9) and (G10) proves the following precise intermediate
obstruction. No estimate

\[
D(Q,\mathcal E_d)-1\le C_{d,r,R}s(Q,\mathcal W)^\beta
\]

can hold for all bodies and all coordinate line/plane decompositions
with \(rB\subseteq Q\subseteq RB\) when \(\beta>1/(d-1)\). One may take
\(r=2/3\), \(R=\sqrt d\), and the same fixed partition throughout:
\(K_t\) contains \((1-t/d)[-1,1]^d\), so these bounds hold for
\(0<t\le1\). This obstruction does \textbf{not} prove that the present
geometric \(1/d\) exponent is sharp as a function of \(s\). The exact
normal measurement has a different scale from relative missing volume.

\subsubsection{\texorpdfstring{4. A rigorous \(1/(d-1)\) recovery
theorem for this full weighted
family}{4. A rigorous 1/(d-1) recovery theorem for this full weighted family}}\label{a-rigorous-1d-1-recovery-theorem-for-this-full-weighted-family}

Assume in addition \(t\le1/2\). Choose a coordinate block attaining the
maximum projected norm of \(a\). At least one index \(k\) is outside
that block because every block has dimension at most two and \(d\ge3\).
Since all coefficients are at most one,

\[
A_*\le a_k\le\operatorname{dist}(a,\mathcal W).
\]

Using \(V\le2^d\) in (G6) gives

\[
s\ge\frac{t^{d-1}}{d!2^{d-1}}.
\]

Since \(S\ge1\), (G8) yields

\[
D(Q_{a,t},\mathcal E_d)-1
\le\frac{t}{S-t}\le2t
\le\boxed{2\bigl(d!2^{d-1}s\bigr)^{1/(d-1)}}.             \tag{G11}
\]

This restricted theorem is uniform over arbitrarily small coefficients

\(a_i\), subject only to \(t\le\min_i a_i\): nearly coordinate normals
do not destroy the improved exponent for weighted corner cuts. Together
with (G9)--(G10), it identifies \(1/(d-1)\) as the optimal geometry
exponent within this family. It does not silently assert that the
general body is a polytope or belongs to this family.

\subsubsection{5. An obstruction to an attractive extra-factor
argument}\label{an-obstruction-to-an-attractive-extra-factor-argument}

It is tempting to replace \(v\le C s\) by \(v\le C\tau s\), with
\(\tau=1-\lambda\). That stronger intermediate assertion is false, even
with fixed radius and inradius bounds and in the explicit family above.

Take \(a=(1,\varepsilon,\ldots,\varepsilon)\), \(t=\varepsilon\),
\(0<\varepsilon\le1/2\), and all coordinate lines as blocks. Then

\[
S=1+(d-1)\varepsilon,\quad
V=2^d-\frac{2\varepsilon}{d!},\quad
\tau=\frac{\varepsilon}{1+(d-1)\varepsilon},
\]

\[
s=\frac{2\sqrt{d-1}\,\varepsilon}{d(d-1)!V},\qquad
v=\frac{2\varepsilon}{d(d-1)!V},\qquad
\frac{v}{\tau s}
=\frac{1+(d-1)\varepsilon}{\sqrt{d-1}\,\varepsilon}
\longrightarrow\infty.                                  \tag{G12}
\]

In particular no constant depending only on \(d\), \(r\), and \(R\) can
justify this factor. Here \(Q\subseteq\sqrt d B\) and, by (G8),
\(Q\supseteq(1-\varepsilon)[-1,1]^d\supseteq\tfrac12 B\).

Even the pointwise proposed estimate

\[
h_P(X)-1\le C\tau\operatorname{dist}(X,\mathcal W)
\]

fails at both new facet samples: the left-hand side is
\(\varepsilon/H\), whereas the distance is \(\sqrt{d-1}\varepsilon/H\).
The ratio before inserting \(\tau\) is fixed. The small facets
responsible for a missing corner can instead be macroscopic nearly
coordinate facets, with a normal angle of order \(\varepsilon\). This is
an actual convex-body counterexample, not a synthetic law unrelated to
its body.

\subsubsection{6. What this resolves and what
remains}\label{what-this-resolves-and-what-remains}

The cap-to-distance step is sharp \textbf{under its volume-only input}.
The separate note \texttt{geometry-rigidity.md} now proves the general
\(1/(d-1)\) recovery theorem under fixed radius and inradius bounds,
with explicit constants. Its direct proof thickens a two-dimensional
convex section and applies the coarea formula to the off-block surface
integral (G1). Thus (G9)--(G10) identify the sharp geometric exponent
under these stronger inputs. No extra \(\tau\) factor is inserted in
(G2), and the volume-only cap conversion is bypassed.

The half-mass and matching stages may also be coupled directly to the
support surplus \(v\), instead of first bounding \(s\). Equations (G7)
and (G12) show why such a coupling must preserve information about the
normal direction and cannot depend only on global containment error.

The all-product distance lower bound (G10) is imported only from the
completed audited theorem; its proof handles arbitrary oblique planar
factor products and does not assume uniqueness of the decomposition.

\subsubsection{7. Exact finite
verification}\label{exact-finite-verification}

The adjacent script \texttt{check\_geometry\_next.py} uses rational
arithmetic to verify facet normalization, the exact \(v\) identity,
squared \(s\) identities, and the near-coordinate counterexample for
finitely many dimensions and rational weights. It is supplementary
verification, not a proof of any universal statement. It needs only
Python's standard library, uses explicit exceptions, and retains checks
under \texttt{python\ -O}.

The corner-cut examples distinguish missing volume of order \(t^d\) from
off-block normal mass of order \(t^{d-1}\). The geometric \(1/(d-1)\)
estimate is established separately; the extra containment factor
proposed for mixed-volume recovery is contradicted by the displayed
family.

\newpage

\section{VII Slanted actual cone law
obstruction}\label{vii-slanted-actual-cone-law-obstruction}

\subsection{Coupled upper-end stability: a cone-law obstruction to a
linear volume
gate}\label{coupled-upper-end-stability-a-cone-law-obstruction-to-a-linear-volume-gate}

This note is separate from the completed upper-end stability supplement.
It contains proved gate obstructions and one explicitly unproved route.
No claim of an improved global exponent is made here. The separate
matching and geometry notes improve the matching and containment gates,
respectively.

\subsubsection{1. Precise obstruction under the selected-basis
intermediate
hypotheses}\label{precise-obstruction-under-the-selected-basis-intermediate-hypotheses}

Let \texttt{0\textless{}e\textless{}=1/2} and put

\[
 Q_e=[-1,1]^3\cap\{|x_1+e x_2|\le1,\ |x_2+e x_3|\le1\}.
 \tag{1.1}
\]

The body is centrally symmetric and full-dimensional, with
\texttt{(1-e){[}-1,1{]}\^{}3\ subset\ Q\_e\ subset\ {[}-1,1{]}\^{}3}.
Its cone law has five direction pairs

\[
 \pm z_1=\pm e_1,\quad \pm z_2=\pm e_2,\quad \pm z_3=\pm e_3,
 \quad \pm z_4=\pm(1,e,0),\quad \pm z_5=\pm(0,1,e).
 \tag{1.2}
\]

Write \texttt{p\_i} for the probability of the entire pair, with uniform
signs inside each pair. Define

\[
\begin{aligned}
 V_e&=8-4e+e^2-e^3/3,\\
 C_1=C_4&=2-e/2,\qquad C_2=2-e,\\
 C_3&=4-3e+e^2-e^3/2,\qquad C_5=2-e+e^2/2,\\
 p_i&=2C_i/(3V_e).
\end{aligned} \tag{1.3}
\]

In particular \texttt{sum\ p\_i=1},
\texttt{p\_4,p\_5\textgreater{}=1/8}, and

\[
 (p_1,p_2,p_3,p_4,p_5)\longrightarrow(1/6,1/6,1/3,1/6,1/6).
 \tag{1.4}
\]

Let \texttt{phi} be the fourth-largest absolute cofactor of a
four-sample tuple, and let
\texttt{A=E\textbar{}det(X\_1,X\_2,X\_3)\textbar{}}. Exact enumeration
gives

\[
 \mathbb E\phi=24p_1p_3p_4p_5e^2\sim e^2/27,
 \tag{1.5}
\]

\[
\begin{aligned}
 A=6[&p_1p_2p_3+e p_1p_2p_5+e p_1p_3p_4+p_1p_3p_5\\
     &+e^2p_1p_4p_5+p_2p_3p_4+e p_2p_4p_5+p_3p_4p_5]
      \longrightarrow2/9,
\end{aligned}\tag{1.6}
\]

and the scalar upper-end deficit is

\[
 \delta_e=\tfrac12-a(Q_e)
     =\frac{3p_1p_3p_4p_5e^2}{A}\sim e^2/48.
 \tag{1.7}
\]

For the basis \texttt{U=I}, the simultaneous section cost used in the
matching lemma is exactly

\[
 Q(U)=\int\phi(e_1,e_2,e_3,x)\,d\nu(x)
   +\sum_{j=1}^3\iint\phi(e_1,\ldots,\widehat e_j,\ldots,e_3,x,y)
                      \,d\nu(x)d\nu(y)
       =2p_4p_5e^2\sim e^2/18.
 \tag{1.8}
\]

Let \texttt{W} be any coordinate partition into lines and planes, and
let \texttt{P\_W=product\_j\ proj\_(W\_j)\ Q\_e}. All four possible
partitions satisfy

\[
\begin{aligned}
 v_W&:=\mathbb E_{\nu_{Q_e}}[h_{P_W}(X)-1]\\
 &=\mathbb E_{\nu_{Q_e}}\operatorname{dist}(X,\bigcup_j W_j)
 \in\{p_4e,p_5e,(p_4+p_5)e\},\\
 \min_W v_W&=e\min(p_4,p_5)\ge e/8.
\end{aligned}\tag{1.9}
\]

Thus neither a coordinate-block mean-distance estimate nor a
coordinate-block mixed-volume-surplus estimate with exponent greater
than \texttt{1/2} in \texttt{Q(U)} can hold under these hypotheses. In
particular, the tempting replacement

\[
 \min_W v_W\le C\mathbb E\phi
 \quad\hbox{or}\quad
 \min_W v_W\le C\delta_e
 \tag{1.10}
\]

is false for this uniformly nondegenerate family. These particular
support calculations concern the prescribed basis and its coordinate
partitions. Section 3 gives separate proofs for arbitrary block choices
and for distance to the entire affine equality class; neither follows
merely from (1.9).

The unit-ball hypothesis of the matching lemma is obtained without
changing the obstruction. Put \texttt{c=sqrt(1+e\^{}2)}, replace
\texttt{Q\_e} by \texttt{c\ Q\_e}, and choose the basis with columns
\texttt{e\_i/c}. Every cone-law sample then has norm at most one, the
basis determinant is \texttt{c\^{}(-3)\textgreater{}=\ (5/4)\^{}(-3/2)},
its coordinate law remains exactly (1.2), and its transformed primal
body remains \texttt{Q\_e}. The section cost and \texttt{E\ phi} are
simply divided by \texttt{c\^{}3}. Its determinant expectation remains
bounded below by a positive constant. Consequently the example meets the
bounded-law, determinant, simultaneous-section, and controlled-basis
hypotheses with constants independent of \texttt{e}.

\subsubsection{2. Proof of the formulas}\label{proof-of-the-formulas}

Fix \texttt{u=x\_2}. The allowed \texttt{x\_1} interval has length
\texttt{2-e\textbar{}u\textbar{}}. The allowed \texttt{x\_3} interval
has length \texttt{2} if \texttt{\textbar{}u\textbar{}\textless{}=1-e},
and length \texttt{1+(1-\textbar{}u\textbar{})/e} otherwise. Therefore

\[
 |Q_e|=2\int_0^1(2-eu)
     \begin{cases}2,&u\le1-e,\\1+(1-u)/e,&u>1-e\end{cases}
          du=V_e.
\]

On the positive \texttt{x\_1=1} facet one has
\texttt{-1\textless{}=u\textless{}=0}, giving area \texttt{2-e/2}. On
the positive \texttt{x\_2=1} facet the allowed \texttt{x\_1} and
\texttt{x\_3} lengths are \texttt{2-e} and \texttt{1}, respectively. On
the positive \texttt{x\_3=1} facet one integrates
\texttt{2-e\textbar{}u\textbar{}} over
\texttt{-1\textless{}=u\textless{}=1-e}, giving \texttt{C\_3}. These
coordinate facets have support number one.

The positive new facet \texttt{x\_1+e\ x\_2=1} is parametrized by
\texttt{x\_1=1-eu}, \texttt{0\textless{}=u\textless{}=1}, and its
allowed \texttt{x\_3} interval. Its area-normal vector is
\texttt{C\_4(1,e,0)}; its area times support number is
\texttt{C\_4=2-e/2}. The other positive new facet has
\texttt{x\_2=1-e\ x\_3}, \texttt{0\textless{}=x\_3\textless{}=1}, and
\texttt{x\_1} length \texttt{2-e+e\^{}2x\_3}. Its area-normal vector is
\texttt{C\_5(0,1,e)}, with area times support number \texttt{C\_5}.
These calculations prove the cone masses in (1.3). They also give
\texttt{sum\ C\_i=3V\_e/2}, as required by the divergence theorem.

For four samples to have four nonzero cofactors, their four distinct
direction types must be \texttt{\{1,3,4,5\}}. Every other type set has a
zero cofactor, and any repeated type also produces at most three nonzero
cofactors. For the exceptional type set the cofactor magnitudes are

\[
 (1,1,e,e^2).
\]

Its fourth absolute cofactor is \texttt{e\^{}2}. Its unnormalized
balanced Rademacher defect is exactly

\[
 \frac{2+e+e^2}{2}
 -\mathbb E|\epsilon_1+\epsilon_2+e\epsilon_3+e^2\epsilon_4|
 =\frac{e^2}{2}.
\]

Indeed, equal signs on the first two terms contribute mean absolute
value \texttt{2}, and opposite signs contribute mean absolute value
\texttt{e}; both cases have probability one half. All other tuples have
zero defect by the balanced three-support equality. The probability of
the exceptional four-type set is \texttt{24p\_1p\_3p\_4p\_5}. Equations
(1.5) and (1.7) follow, using \texttt{Delta=4A\ delta}. Enumerating the
ten three-type sets gives (1.6).

For (1.8), the one-sample section vanishes because every sample is at
most two-sparse. Among the two-sample sections, only the one omitting
\texttt{e\_2} can have four nonzero cofactors; its two additional types
must be \texttt{4,5}, in either order. Its fourth cofactor is
\texttt{e\^{}2}, giving the formula.

Every coordinate projection has both axis support numbers equal to one.
The projection onto the \texttt{12} plane preserves the first new
inequality, and the projection onto the \texttt{23} plane preserves the
second. Thus the \texttt{12\textbar{}3} product has support surplus zero
at \texttt{z\_4} and \texttt{e} at \texttt{z\_5}; the
\texttt{1\textbar{}23} product has the opposite behavior. At the
\texttt{13\textbar{}2} product and the all-line product both new normals
have surplus \texttt{e}. All axis normals have surplus zero. The
Euclidean distances to the block unions are precisely the same: an
unhandled pair has a tail of length \texttt{e} and all handled pairs
have distance zero. This proves (1.9).

\subsubsection{3. What this says about a coupled
proof}\label{what-this-says-about-a-coupled-proof}

\paragraph{Sharp square-root matching even with arbitrary block
choices}\label{sharp-square-root-matching-even-with-arbitrary-block-choices}

The mean-distance obstruction is stronger than the fixed-coordinate
statement. It persists if the line and plane blocks, and the sampled
basis, can be chosen freely. The following elementary finite-set lemma
makes this precise.

\textbf{Finite-set lemma.} Let \texttt{d\textgreater{}=3},
\texttt{0\textless{}e\textless{}=1/2}, and let
\texttt{R\^{}d=direct\_sum\ W\_j}, with \texttt{dim\ W\_j} equal to one
or two. For

\[
 \mathcal Z_e=\{e_1,\ldots,e_d,e_1+e e_2,e_2+e e_3\}
\]

one has

\[
 \max_{z\in\mathcal Z_e}\operatorname{dist}(z,\bigcup_jW_j)
       \ge\frac{e}{6\sqrt d}. \tag{3.2}
\]

\textbf{Proof.} Suppose the maximum is
\texttt{r\textless{}e/(6\ sqrt\ d)}. Assign each axis vector to a
nearest block and choose \texttt{w\_i} in that block with
\texttt{\textbar{}w\_i-e\_i\textbar{}\textless{}=r}. Let \texttt{T} have
columns \texttt{w\_i}. Then
\texttt{\textbar{}\textbar{}T-I\textbar{}\textbar{}\textless{}=sqrt\ d\ r=:eta\textless{}1/12},
so \texttt{T} is invertible. The number of assigned axis indices in each
block is at most its dimension. Summing the numbers and dimensions gives
\texttt{d} on both sides; therefore the assigned index sets
\texttt{J\_j} partition the axes, have sizes one or two, and
\texttt{W\_j=T\ span\{e\_i:i\ in\ J\_j\}}.

For \texttt{y} in \texttt{W\_j},

\[
 |\operatorname{proj}_{J_j^c}y|
       \le\frac\eta{1-\eta}|y|.
\]

Take \texttt{y} to be the orthogonal projection of an arbitrary
\texttt{z} onto \texttt{W\_j}. Since
\texttt{\textbar{}y\textbar{}\textless{}=\textbar{}z\textbar{}}, the
triangle inequality gives

\[
 \operatorname{dist}(z,W_j)
 \ge\operatorname{dist}(z,\operatorname{span}J_j)
          -\frac\eta{1-\eta}|z|.
\]

The two edges \texttt{12} and \texttt{23} cannot both be contained in
blocks of a partition whose block sizes are at most two. For the
corresponding uncontained edge vector \texttt{z}, its distance to every
coordinate block is at least \texttt{e}. Since
\texttt{\textbar{}z\textbar{}\textless{}=sqrt(5/4)}, its distance to
every actual block is at least
\texttt{e-3\ sqrt(d)\ r\textgreater{}e/2\textgreater{}r}, a
contradiction. This proves (3.2).

For every dimension, take the actual cone body

\[
 K_{e,d}=Q_e\times[-1,1]^{d-3}. \tag{3.3}
\]

Its cone law has the five three-dimensional direction-pair masses
\texttt{(3/d)p\_i}, and mass \texttt{1/d} on each additional
coordinate-axis pair. All \texttt{d+2} masses are at least
\texttt{3/(8d)}. Applying (3.2) therefore yields

\[
 \inf_{(W_j)}\mathbb E_{\nu_{K_{e,d}}}
              \operatorname{dist}(X,\bigcup_jW_j)
       \ge\frac{e}{16d^{3/2}}. \tag{3.4}
\]

The determinant and fourth-cofactor expectations are exactly

\[
 A_d=\frac{d!\,27}{6d^d}A,
 \qquad
 \mathbb E\phi_d=\frac{(d+1)!\,81}{d^{d+1}}
                 p_1p_3p_4p_5e^2
       \sim\frac{(d+1)!}{8d^{d+1}}e^2. \tag{3.5}
\]

For the first formula, a determinant needs one sample on each additional
coordinate axis and three independent samples in the first
three-dimensional subspace. For a positive fourth cofactor, a
\texttt{(d+1)}-tuple needs one sample on every additional axis and the
four exceptional types \texttt{\{1,3,4,5\}} in that subspace. Its four
nonzero cofactors have magnitudes \texttt{(1,1,e,e\^{}2)}, with all
other cofactors zero. This proves (3.5) by ordered-type counting. The
scalar deficit is

\[
 \delta(K_{e,d})=(3/d)\delta_e\sim e^2/(16d), \tag{3.6}
\]

either by the dimension-weighted product identity or by the same exact
defect count. The section cost at the axis basis is
\texttt{18\ p\_4\ p\_5\ e\^{}2/d\^{}2}.

Scaling by \texttt{c=sqrt(1+e\^{}2)} puts the entire cone law in the
unit ball. Every basis matrix formed from its samples has
\texttt{\textbar{}\textbar{}U\textbar{}\textbar{}\textless{}=sqrt\ d}.
If its coordinate blocks are \texttt{W\textquotesingle{}\_j}, their
original blocks are \texttt{U\ W\textquotesingle{}\_j}, and

\[
 \operatorname{dist}(U^{-1}X,\bigcup_j W'_j)
 \ge \frac{1}{\sqrt d}
           \operatorname{dist}(X,\bigcup_j UW'_j).
\]

Consequently (3.4) gives a coordinate mean lower bound
\texttt{e/(16\ c\ d\^{}2)} for \textbf{every} invertible sampled basis
and every line/plane partition. The bounded-law determinant mass has an
explicit uniform lower bound: using \texttt{p\_i\textgreater{}=1/8} and
\texttt{A\textgreater{}=6p\_1p\_2p\_3\textgreater{}=3/256}, one obtains

\[
 A_d/c^d\ge
   \frac{27d!}{512d^d(5/4)^{d/2}}. \tag{3.7}
\]

It follows that, in every fixed dimension \texttt{d\textgreater{}=3}, no
exponent greater than \texttt{1/2} can replace the square-root exponent
in a mean-distance matching lemma under its actual hypotheses, even when
the law is required to be the cone law of a uniformly nondegenerate
symmetric body and the block decomposition can be chosen freely.

\paragraph{The discrepancy is also linear in distance to the entire
product
class}\label{the-discrepancy-is-also-linear-in-distance-to-the-entire-product-class}

The same family satisfies a full-class geometric estimate:

\[
 \boxed{\frac{e}{120\sqrt d}
       \le D(K_{e,d},\mathcal E_d)-1
       \le\frac{e}{1-e}\le2e.} \tag{3.8}
\]

This does not improve the existing \texttt{1/d} obstruction, since here
the scalar deficit is of order \texttt{e\^{}2}. It shows, however, that
the matching square-root loss can reflect an actual first-order product
discrepancy, rather than only a badly chosen coordinate presentation.

We use a general elementary flat-facet observation. Suppose a
full-dimensional body \texttt{K} has a facet with unit normal
\texttt{n}, support number \texttt{h\textgreater{}0}, and a Euclidean
tangent disk of radius \texttt{rho} contained in the facet. If
\texttt{K\ subset\ E\ subset\ (1+s)K}, the upper boundary of \texttt{E},
over that disk and in direction \texttt{n}, is the graph of a finite
concave function with height in \texttt{{[}0,sh{]}} above the facet
plane. At a differentiability point in the concentric disk of radius
\texttt{rho/2}, its gradient has norm at most \texttt{2sh/rho}: restrict
the concave function to each tangent line, and compare its derivative
with secants to points at distance \texttt{rho/2} on either side. Such
differentiability points exist by local Lipschitz continuity of finite
concave functions. Their outward unit normal \texttt{n\_E} satisfies

\[
 \operatorname{dist}(n,\operatorname{span}n_E)
       \le\frac{2sh}{\rho}. \tag{3.9}
\]

If \texttt{E} is an affine product of line and plane factors, every
regular outward normal of \texttt{E} lies in one of its dual blocks, a
direct-sum partition into subspaces of dimensions one and two. Indeed,
if two primal factors were on their boundaries at a point, their nonzero
normal cones would give at least two independent normal rays; that point
could not be regular. At a regular point exactly one factor is on its
boundary, and its normal annihilates all the other primal blocks.

For \texttt{Q\_e}, tangent disks of radius \texttt{rho=1/10} lie in its
five positive facets, with centers

\[
 (1,-1/2,0),\quad(0,1,-1/2),\quad(0,0,1),\quad
 (1-e/2,1/2,0),\quad(0,1-e/2,1/2).
\]

For the two new facets the disk is taken in their actual tangent planes.
All disks satisfy the remaining inequalities with room to spare for
\texttt{e\textless{}=1/2}. In \texttt{K\_\{e,d\}}, use the same centers
with zero extra coordinates; the radius-\texttt{1/10} disks in the full
\texttt{(d-1)}-dimensional tangent planes remain inside the
corresponding facets. Each additional coordinate facet also contains
such a disk centered at its coordinate-axis endpoint. All facets in
question have the scaled normals in \texttt{Z\_e}, namely
\texttt{z=n/h}.

For any \texttt{E\ in\ E\_d} with
\texttt{K\_\{e,d\}\ subset\ E\ subset\ (1+s)K\_\{e,d\}}, (3.9) therefore
puts every \texttt{z\ in\ Z\_e} within

\[
 (1/h)(2sh/\rho)=20s
\]

of the union of \texttt{E}'s dual line/plane blocks. Equation (3.2)
forces \texttt{20s\textgreater{}=e/(6\ sqrt\ d)}. Every Banach--Mazur
sandwich can be rescaled to this form, and taking its infimum proves the
lower bound in (3.8), without any attainment assumption. Finally
\texttt{(1-e){[}-1,1{]}\^{}d\ subset\ K\_\{e,d\}\ subset\ {[}-1,1{]}\^{}d}
gives the upper bound.

\paragraph{A joint geometric gate remains
necessary}\label{a-joint-geometric-gate-remains-necessary}

The family has deficit of order \texttt{e\^{}2} and a ridge-scale
geometric discrepancy of order \texttt{e}. It does not threaten the
conjectured optimal dimension-three power \texttt{1/3}, since \texttt{e}
is smaller than a constant times \texttt{delta\^{}(1/3)}. Instead it
shows why it is unsafe to strengthen the final volume gate by inserting
a square into the mean-distance estimate. The small determinant factor
\texttt{e\^{}2} is compensated by geometry spread over a ridge, whose
cap-volume scale has exponent two rather than three.

By contrast, the three-coordinate corner truncation has \texttt{E\ phi}
of order \texttt{t\^{}2}, \texttt{delta} of order \texttt{t\^{}3}, and
support surplus of order \texttt{t\^{}3}. In that case the half-mass
gate and a support-Lipschitz estimate lose a factor which the actual
support cost already contains. The two families show different gates
becoming active. The separate extreme bounds cannot be assumed to
saturate on one body.

There is a useful exact weighted estimate. For the normalized cofactor
tuple, let \texttt{g=1/2-t\_1} and \texttt{phi=M\ t\_4} on
\texttt{M\textgreater{}0}, and define \texttt{g=phi=0} on \texttt{M=0}.
Since \texttt{t\_4\textless{}=1/4} and \texttt{g\textless{}=1/2} on
\texttt{M\textgreater{}0}, Lemma 3.1 of the completed proof implies

\[
 t_4g\le\min(t_4/2,g/4)\le\mathscr D(t),
 \qquad \boxed{\mathbb E[\phi g]\le NA\delta.}
 \tag{3.1}
\]

This carries no power loss, but it cannot replace the omitted-basis
matching cost by itself. Matching consistency must still be proved using
a geometric cost which treats a shallow long circuit and two overlapping
nearly-axis pair directions differently.

An unproved route toward the optimal \texttt{1/d} is to produce a
geometric witness at a product-containment defect \texttt{t}, classify
the support involved by its active block dimension, and show that the
cone-law masses times the circuit defect are at least
\texttt{c\_d\ t\^{}r} for some \texttt{r\textless{}=d}. A proof must
account for determinant factors, the normal-angle factors in localized
facet mass, and non-polyhedral contact patches together. Merely
optimizing the independent mean estimates does not establish such a
statement. No complete lemma of this form is proved in this note.

\subsubsection{4. Exact finite
verification}\label{exact-finite-verification-1}

\texttt{check\_global\_slanted.py} replays the cone-mass identity, the
determinant expectation, all 625 ordered direction-type four-tuples with
their 16 sign averages, all simultaneous two-sample sections, and the
four block costs, at eight rational values of \texttt{e}. It also checks
the product determinant, fourth-cofactor and scalar-deficit formulas in
32 cases across dimensions three through six. It uses exact rational
arithmetic and exceptions rather than Python assertions. The finite
checks support the displayed identities; the all-parameter proof is the
algebra above.

The finite-set lemma, facet disks, and flat-facet geometric argument in
Section 3 also received a separate analytical audit, with no gap found.
This is model-assisted mathematical review, not external peer review or
formal certification; the exact tests are not used to establish these
geometric arguments.

\end{document}
